% Research note: turns of a T-mesh, obstructions to closing edge loops (the
% theory paper; the applied companion is docs/research/closing-loops). The Lean
% project that checks its algebraic cores is in lean/, beside this file.
% Build: `make paper NOTE=turns-of-a-t-mesh` from the root of the repository,
% which first writes version.tex (the commit described, \enginecommit);
% without it, or by hand: pdflatex paper && bibtex paper && pdflatex paper && pdflatex paper
\documentclass[11pt,a4paper]{article}

\usepackage[T1]{fontenc}
\usepackage[utf8]{inputenc}
\usepackage{lmodern}
\usepackage[margin=2.5cm]{geometry}
\usepackage{amsmath,amssymb,amsthm}
\usepackage{booktabs}
\usepackage{array}
\usepackage[font=small]{caption}
\usepackage{tikz}
\usetikzlibrary{arrows.meta,decorations.pathreplacing,patterns}
\usepackage[numbers,sort&compress]{natbib}
\usepackage[hidelinks]{hyperref}
\setlength{\emergencystretch}{2em}
\renewcommand{\topfraction}{0.9}
\renewcommand{\floatpagefraction}{0.8}
% The references in a smaller type.
\renewcommand{\bibfont}{\small}
\setlength{\bibsep}{2pt plus 1pt}

\newtheorem{theorem}{Theorem}
\newtheorem{definition}{Definition}
\newtheorem{proposition}{Proposition}
\newtheorem{lemma}{Lemma}
\newtheorem{corollary}{Corollary}
\newtheorem{remark}{Remark}
\newtheorem{conjecture}{Conjecture}
\newtheorem{question}{Question}

\newcommand{\Z}{\mathbb{Z}}
\newcommand{\R}{\mathbb{R}}
\newcommand{\shift}{s}
\newcommand{\RC}{\textup{\textsc{Ring-Closing}}}

% The commit of the Lean proofs described, written by `make paper` into
% version.tex (ignored by git).
\IfFileExists{version.tex}{\input{version.tex}}{}
\providecommand{\enginecommit}{unknown}

\title{Turns of a T-Mesh:\\
Obstructions to Closing Edge Loops\\[0.5em]
\large Research note --- draft, not peer reviewed}
\author{Marc Gardent%
\thanks{This note is distributed under the Creative Commons Attribution
4.0 International licence (CC BY 4.0). The code of openclare and the Lean
proofs are under the GPL-3.0-or-later of the repository. The note describes
the Lean proofs at commit \texttt{\enginecommit} of
\url{https://github.com/marcgardent/openclare}.
A French translation is in \texttt{paper.fr.pdf}, beside this note.}\\
\normalsize Independent researcher\\
\normalsize ORCID: \href{https://orcid.org/0009-0000-6466-3098}{0009-0000-6466-3098}\\
\normalsize Contact: through the repository,
\url{https://github.com/marcgardent/openclare}}
\date{October 2026}

\begin{document}
\maketitle

\begin{abstract}
In the quantization of a T-mesh, an edge loop closes into a ring when a
\emph{turn}, a straight walk across the patches that returns to its first
side, has zero \emph{shift}, an integer linear form in the arc lengths;
otherwise it winds as a helix. We study why turns cannot all be closed.
Reading the shift as translational holonomy of the flat metric, we show that
for a simple turn it is the across component of the Burgers vector of the
singular vertices on one side; that a 3--5 pair or an exchange of two
vertices, joined by a staircase of arcs, forces one of two turns to be
shifted, and that a separatrix winding once round a turn and stopping on
another separatrix of its vertex forces that turn alone to be shifted; and
that on a
closed sphere the shift and the length of a turn have the same parity. Fox calculus extends the Burgers sum to
any combination of turns, and \emph{combs}, local dislocations carried by
single arcs, write every certificate of a conflict found on a human body. The
conflicts bound from below the number of turns that every quantization within
given bounds leaves shifted: on that body, within the bounds of separation of
a remeshing engine, exactly one to nine of the turns it holds, where it leaves
13 to 123 shifted. Choosing the turns to close is NP-hard for bidirected
balances with $0$--$1$ bounds, the algebraic form of quantization; whether
actual T-meshes give hard instances is open. The algebraic cores
of the proofs are checked in Lean; the topology of the surface, that the walks
and regions are those of a T-mesh, is a hypothesis there.
\end{abstract}

\section{Introduction}

Automatic quad remeshing often produces edge loops that wind round a part of
the shape as helices instead of closing into
rings~\cite{bommes2011structure}. Whether a loop closes is decided when a
quantizer assigns whole lengths to the arcs of a T-mesh traced in a seamless
parametrization~\cite{campen2015qgp,lyon2021constrained,heistermann2023bimdf}:
a loop follows a \emph{turn}, a straight walk across the patches that returns
to its first side, and closes when the \emph{shift} of the turn, an integer
linear form in the arc lengths, vanishes. A companion
paper~\cite{gardent2026loops} adds these equations to a quantizer; on real
models they cannot all hold together with the bounds that keep singular
vertices apart. This paper asks why, and what follows.

Sections~\ref{sec:setting} to~\ref{sec:beyond} go from the holonomy of one
turn to conflicts between turns, any combination of turns, and borders and
handles; Sections~\ref{sec:forced} and~\ref{sec:choosing} bound the turns
every quantization leaves shifted and show that choosing them is NP-hard;
Sections~\ref{sec:lean} and~\ref{sec:open} say what is machine-checked and
what is open. The ingredients are classical (holonomy of flat cone metrics, Burgers vectors,
the parity of quadrangulations, Fox calculus, covers of infeasible
subsystems); their statements for turns, the lower bound on helices and the
hardness for the instances of quantization are ours, as far as we found.
Statements are \emph{proved} when they come with a proof, \emph{measured}
when they come from exact arithmetic or solver output on the systems the
remesher openclare exports, and \emph{conjectured} otherwise.

This work began as a programming project, openclare, an open-source quad
remesher, and became research when engineering alone could no longer tell
which results could be guaranteed and which computations saved. The results
here answer from the side of what cannot be done: a repair that keeps
the length of a turn of odd shift cannot close it on a closed sphere; a turn
whose shift is a sign-definite form with an arc held open stays a helix and
need not be tried; the turns every quantization leaves shifted measure what a
heuristic can still gain; and the hardness of choosing them is why an engine
settles for a heuristic.

\section{Related work}\label{sec:related}

The valences of a closed quad mesh of genus $g$ satisfy $\sum_v (4 - k(v)) = 8(1-g)$,
which says how many singular vertices are needed, not where. A closed quad mesh induces a quartic differential
whose zeros and poles are its singular vertices; Lei et
al.~\cite{lei2020quartic} list as open the guarantee that all its
trajectories are finite, and a trajectory that is not finite is an edge loop
that does not close. A quadratic
differential whose trajectories all close is
Jenkins--Strebel~\cite{strebel1984quadratic}, a union of flat cylinders; a
helix is a cylinder glued back with a shift of whole rows. Chen et
al.~\cite{chen2019metric} characterize the flat metrics of quad meshes,
assuming that the lines of the field are finite loops. Izmestiev et al.~\cite{izmestiev2013torus} prove non-existence results for quad meshes,
among them the torus with one vertex of valence three and one of valence
five~\cite{barnette1971toroidal}, with the holonomy of the developing map,
whose translation, when the rotation is trivial, is the Burgers vector of
crystallography; a pair of opposite disclinations is a
dislocation~\cite{kleman2008disclinations}. The periods of a flat cone metric
are local coordinates of it~\cite{thurston1998shapes,veech1993flat}. Bommes et
al.~\cite{bommes2011structure} remove helices from finished meshes; stripe
parameterizations of tubes~\cite{kalberer2011stripe} round the gap round a
tube, our shift, to a whole number. We found no lower bound on the number of helices for a given
singularity configuration; ours is a minimum cover of irreducible infeasible
subsystems~\cite{parker1996iis,amaldi2003maxfs}.

\section{Turns, shifts and holonomy}\label{sec:setting}

A T-mesh is a partition of the surface into rectangles, the \emph{patches},
whose sides are made of \emph{arcs}; nodes where a trace stops on another are
T-junctions, and the singular vertices of the field are nodes. Each patch $P$
has local coordinates with its first corner at the origin and its sides along
the axes. A quantization assigns every arc $a$ a length $\ell_a \in \Z_{\ge
0}$ such that opposite sides of every patch have equal length, the
\emph{balances} $B\ell = 0$; the quad mesh is then the integer grid of each
patch.

\begin{definition}[Turn]
A \emph{turn} $\gamma$ is a straight walk across patches, from an offset of
its first side that meets no node: it enters a patch through a side, leaves
through the opposite side at the same distance from the corresponding corner,
enters the neighbouring patch across the arc it reaches, and so on, until it
comes back into its first patch through its first side. The rotation of the
composed crossings is then the identity, since the walk re-enters its first
side inward.
\end{definition}

Crossing from patch $P$ to patch $Q$ across an arc composes a quarter-turn
rotation of the local frame with a translation whose components are signed
sums of arc lengths along the rims of $P$ and $Q$.

\begin{proposition}[The shift of a turn is linear]\label{prop:linear}
For a turn $\gamma$ there are integer coefficients $c_a$, independent of the
lengths, such that after a full turn a point of the first side has moved
across the walking direction by
\[
  \shift_\gamma(\ell) = \sum_a c_a\, \ell_a .
\]
For integer lengths $\ell$, let $o$ be a half-integer offset of the first side
whose walk at the lengths $\ell$ crosses the same patches as $\gamma$ (we call
the set of such offsets the \emph{corridor} of $\gamma$ at $\ell$). Then the
dual line of the quad mesh at $o$ comes back at $o + \shift_\gamma(\ell)$: it
closes after one round if and only if $\shift_\gamma(\ell) = 0$. If moreover
$o + j\,\shift_\gamma(\ell)$ lies in the corridor for $j = 1, \dots, t - 1$, it
winds $t$ times round as a helix of pitch $|\shift_\gamma(\ell)|$ in the sense
of~\cite{bommes2011structure}. If the corridor holds no half-integer, no strip
of the mesh follows $\gamma$.
\end{proposition}

\begin{proof}[Sketch]
Each crossing is a map $x \mapsto R_k x + t$ with $R_k$ a rotation by $k$
quarter-turns and $t$ a vector whose entries are sums of $\pm\ell_a$. The
composition along $\gamma$ has trivial rotation, hence is a translation
$T_\gamma(\ell)$, linear in $\ell$ with integer coefficients; its component
across the walking direction is $\shift_\gamma(\ell)$. Singularities are
nodes, so the quad mesh is a regular grid along the walk, and a dual line at
distance $o$ from the first corner that crosses the same patches comes back at
$o + \shift_\gamma(\ell)$. The corridor, two inequalities per crossing, may
contain $o$ but not $o + \shift_\gamma(\ell)$; the next round then follows
other patches.
\end{proof}

In the frame of the first side, with $e_1$ the walking direction and $e_2$
the direction of increasing offsets, $T_\gamma(\ell) = -\alpha_\gamma(\ell)\,
e_1 + \shift_\gamma(\ell)\, e_2$: it brings the end of the walk back onto the
first side. $\alpha_\gamma(\ell)$ is the \emph{length} of the turn, which at whole lengths is the number of
quads of the strip in one turn. At the real lengths of the parametrization the
shift is the \emph{real shift}, and its share of the length the \emph{lean}.
At whole balanced lengths with positive patch sides every strip of the quad
mesh is such a walk~\cite[Prop.~2]{gardent2026loops}, and every shift is
whole; closing a ring asks it to be zero.

\paragraph{Flat metric and holonomy.}
Unless stated otherwise the surface $\Sigma$ is closed, and \emph{genus zero}
means a closed sphere. Balanced lengths $\ell > 0$ make each patch a
Euclidean rectangle, and $\Sigma$ a flat surface with cone points at the
singular vertices: a vertex of valence $k$ has angle $k\pi/2$ and curvature
$(4-k)\pi/2$, and T-junctions are not cone points. Let $\Sigma^\circ$ be $\Sigma$ without its singular vertices,
$x_0$ a point inside a patch, and $R$ the positive quarter-turn. Developing a
loop $\gamma$ of $\Sigma^\circ$ based at $x_0$ patch after patch gives its
\emph{holonomy}, a motion $\hat h(\gamma)\colon x \mapsto R^j x + t$ of the
coordinates of the first patch; it depends only on the homotopy class of
$\gamma$, and $\hat h(\alpha\beta) = \hat h(\alpha) \circ \hat
h(\beta)$~\cite{izmestiev2013torus}. A small loop round a singular vertex $c$
of valence $k$, run with $c$ on its right and reached along a path that
develops it at $p$, has holonomy $x \mapsto R^{4-k}(x - p) + p$. When the
rotation of $\hat h(\gamma)$ is the identity, its translation is the
\emph{Burgers vector} of $\gamma$. For a turn $\gamma$, let $\bar\gamma$ be
the closed curve that follows the walk and, back in its first patch, returns
to its start along a segment parallel to the first side; $\hat
h(\bar\gamma)$ is a translation. It continues the chart along the loop,
while the crossings of Proposition~\ref{prop:linear} develop the last patch
into the first: $\hat h(\bar\gamma)$ is the inverse of $T_\gamma(\ell)$, the
translation by $\alpha_\gamma(\ell)\, e_1 - \shift_\gamma(\ell)\, e_2$, and
$\shift_\gamma(\ell) = -\langle e_2, \hat h(\bar\gamma) \rangle$. Every shift
in this paper is that of Proposition~\ref{prop:linear}, the convention of the
shifts the engine exports and of the figures. Shifts, developed positions and Burgers vectors are
integer linear forms in $\ell$ as long as the paths cross the same patches, so
an identity between them that holds on an open set of balanced positive
lengths holds for all balanced lengths; for whole lengths the nodes develop
onto $\Z^2$. A form $f$ equals $g$ \emph{modulo the balances} if $f - g$ is a
combination of the rows of $B$.

\paragraph{Data.}
The measured statements refer to the T-mesh of the male body of Blender's
Essentials library remeshed by openclare at $10{,}590$ quads, exported with
every turn (branch \texttt{rd-exact}, engine before commit \texttt{c8bf2be})
and studied in the companion paper~\cite[\S7]{gardent2026loops}: $261$
singular vertices ($133$ of valence three, $127$ of valence five, one of
valence two) among $1{,}295$ nodes, $2{,}069$ arcs, $776$ patches, $1{,}552$
balances. Turns are found in eight \emph{ways}, way $m + 2j + 4i$ walking from
the shares $\{0.5\}$ ($m = 0$) or $\{0.2, 0.5, 0.8\}$ ($m = 1$) of each side
and \emph{holding} (trying to close) a turn if its real shift is above $0.05$,
its lean below $0.05$ ($j = 0$) or $0.12$ ($j = 1$) and its length below $2$
($i = 0$) or $6$ ($i = 1$) sides of the T-mesh; the engine keeps way~2. The
ways find $1{,}299$ turns and hold $765$ (summed over the ways). The integer
search holds $28$ to $36$ arcs at length at least one, the \emph{bounds of
separation}; $P_{\mathrm{eng}}$, the
balanced lengths meeting them, is neither a subset nor a superset of the valid
quantizations, which only ask a chain between two singular vertices to keep a
positive length. Identities are checked in floating point unless said exact;
programs were solved with HiGHS 1.8.0~\cite{huangfu2018highs} through SciPy
1.16.3~\cite{virtanen2020scipy}, on an Intel Core i7-7700K (4 cores, 2017,
62\,GB) shared with other jobs.

\section{Turns and singularities}\label{sec:holonomy}

\begin{lemma}[Gauss--Bonnet for turns]\label{lem:gauss}
Let $\gamma$ be a turn whose curve $\bar\gamma$ is simple. If $\bar\gamma$
bounds a disc $D \subset \Sigma$, the singular vertices inside $D$ satisfy
$\sum_{c \in D} (4 - k(c)) = 4$. If the curves of two turns are disjoint and
bound an annulus $A$, then $\sum_{c \in A} (4 - k(c)) = 0$.
\end{lemma}

\begin{proof}
For a region of a flat surface with cone points bounded by straight segments,
Gauss--Bonnet reads $\sum_c \kappa_c + \sum_i \theta_i = 2\pi\chi$, with
$\kappa_c = (4 - k(c))\pi/2$ the curvatures inside and $\theta_i$ the turning
angles of the border~\cite{bommes2013survey}. The curve $\bar\gamma$ turns
only at the ends of its closing segment, by $+\pi/2$ and $-\pi/2$, so the
turning angles sum to zero on each curve, seen from either side; $\chi(D) = 1$
and $\chi(A) = 0$.
\end{proof}

On a closed sphere a simple turn thus splits the singular vertices into two
groups of total index four, each of at least two vertices: a fingertip beyond
a ring holds the equivalent of four vertices of valence three.

\begin{lemma}[One self-crossing]\label{lem:eight}
Let the surface be closed of genus zero, and let the curve $\bar\gamma$ of a
turn cross itself exactly once, transversally, with no singular vertex on it.
It cuts the surface into three discs: two \emph{lobes}, each bounded by one of
the two loops of $\bar\gamma$, and a third region bounded by both. If the
crossing lies on the walk, the totals of $4 - k$ over the singular vertices of
the lobes and the third region are $3$, $3$ and $2$; if it lies on the closing
segment, they are $2$, $2$ and $4$.
\end{lemma}

\begin{proof}
The walk crosses itself at a right angle, so each loop leaves the crossing
$X$ along one passage and comes back along the other, its two rays a
quarter-turn apart. Each loop is simple and bounds a disc containing the
quadrant between its rays, a lobe with a corner of angle $\pi/2$ at $X$; the
lobes take opposite quadrants and the third region the other two, and the
three faces of this one-vertex graph on the sphere are discs. By
Gauss--Bonnet a disc whose border turns by $\theta_i$ holds total index $4 -
\sum_i \theta_i/(\pi/2)$, which gives $3$, $3$ and $2$ before the corners of
the closing segment, of turnings $+\pi/2$ and $-\pi/2$ in succession. If the
crossing lies on the walk, both are on one loop, every region sees both, and
they cancel. Otherwise let the walk start at $A$ along $e_1$, come back at $Z$
and cross the closing segment between them; it crosses along $-e_1$, since
along $e_1$ the turn would have ended there. The loop through $A$ then turns
the same way at $A$ as at $X$, and so does the loop through $Z$ at $Z$ and at $X$ (a
check of the four directions): each lobe sees a second convex corner and holds
$2$, and the third region the rest of the total $8$.
\end{proof}

\begin{proposition}[The shift is set by the singularities]\label{prop:burgers}
Let $\gamma$, $\gamma_1$, $\gamma_2$ be turns with simple curves.
\begin{enumerate}
\item[(a)] Let $\bar\gamma$ bound a disc $D$ holding the singular vertices
  $c_1, \dots, c_m$, and let $x_0$ be a point of $\bar\gamma$, the start of the
  walk for instance. Let $T \subset D$ be a tree whose vertices include $c_1,
  \dots, c_m$ and which meets $\bar\gamma$ only at $x_0$. Cut along $T$, the
  disc $D$ becomes a disc $\Omega$ without cone point, which develops into the
  plane, with the frame $(e_1, e_2)$ at $x_0$. Number the $c_i$ in the order in
  which a walk round $T$ from $x_0$, keeping $T$ on its right, first reaches
  them, and let $p_i$ be the developed position of $c_i$ at that first visit.
  Then $\shift_\gamma(\ell) = \langle e_2, B_D \rangle$ if $D$ lies on the
  left of $\bar\gamma$, and $\shift_\gamma(\ell) = -\langle e_2, B_D
  \rangle$ if it lies on its right, with
  \[
    B_D = \sum_{i=1}^m Q_i\,(I - R_i)\,p_i ,
    \qquad R_i = R^{4 - k(c_i)},
    \qquad Q_i = R_m R_{m-1} \cdots R_{i+1},
  \]
  the translation of $\rho_m \circ \dots \circ \rho_1$, where $\rho_i(x) =
  R_i(x - p_i) + p_i$. With $D$ on its left, $B_D = T_\gamma(\ell) =
  -\alpha_\gamma(\ell)\, e_1 + \shift_\gamma(\ell)\, e_2$; with $D$ on its
  right, $B_D$ is the inverse translation.
\item[(b)] Let $\bar\gamma_1$ and $\bar\gamma_2$ be disjoint and bound an
  annulus $A$, round which the two turns go the same way, and let a path
  $\kappa$ inside $A$ from $\bar\gamma_1$ to $\bar\gamma_2$ carry the walking
  direction of $\gamma_2$ onto that of $\gamma_1$ (the turns are
  \emph{parallel}); measure both shifts along the direction $e_2$ carried by
  $\kappa$. Then
  \(
    \shift_{\gamma_1}(\ell) - \shift_{\gamma_2}(\ell) = -\langle e_2, B_A \rangle,
  \)
  where $B_A$ is the Burgers vector of the border of the disc $A$ cut along
  $\kappa$, given by the same sum over the singular vertices inside $A$,
  developed inside the cut disc.
\item[(c)] In particular, parallel turns bounding an annulus without singular
  vertex have the same shift for every $\ell$: a helix is a property of a
  whole cylinder of parallel turns, not of one of them.
\end{enumerate}
\end{proposition}

\begin{figure}[t]
\centering
% The shift of a turn is set by the singular vertices on one side of it
% (Proposition 2, "The shift is set by the singularities").
% \input inside a figure; needs \usepackage{tikz} and
% \usetikzlibrary{arrows.meta,decorations.pathreplacing}.
%
% A cap D of a T-mesh (a fingertip, say), cut along a tree T into a disc
% Omega that develops into the plane; every node is at a whole position and the
% faint grid is the quad mesh at these lengths. Omega is the band [0,11]x[.,4]
% (it goes on below), the block [3,11]x[4,5] and the top, the L-shaped
% [0,3]x[4,6] + [0,2]x[6,7]. T runs from x_0 up the left side to c_1 and then
% round the border of the top: c_1, c_2, ..., c_6, which are the six singular
% vertices of D. The cut sides are glued in pairs: the right side x = 11 to the
% left side one unit lower, (11,y) ~ (0,y-1); and the top of the block, from
% c_6 = (3,5) rightwards, onto the border of the top, turning round it:
% [3,4] onto [(3,5),(3,6)], [4,5] onto [(3,6),(2,6)], [5,6] onto [(2,6),(2,7)],
% [6,8] onto [(2,7),(0,7)], [8,11] onto [(0,7),(0,4)]. So c_1 = (11,5) = (0,4),
% c_2 = (8,5) = (0,7), c_3 = (6,5) = (2,7), c_4 = (5,5) = (2,6), c_5 = (4,5) =
% (3,6), c_6 = (3,5). Angles: 90 + 90 + 90 at c_1; 90 + 180 at c_2, c_3, c_5;
% 180 + 90 at c_6; 270 + 180 at c_4. Valences 3, 3, 3, 5, 3, 3; total index
% 5 - 1 = 4 (Lemma 1, a disc).
%
% Traces besides the cut: c_1 to the right (stops at (2,4)), c_4 to the left
% (stops on the cut at (0,6) = (9,5)) and downwards (x = 2), c_6 to the left
% (stops at (2,5)) and downwards (x = 3), and the separatrices downwards from
% c_5, c_4, c_3, c_2 at x = 4, 5, 6, 8.
%
% The turn gamma from x_0 = (0, 3.5) crosses x = 2, 3, 4, 5, 6, 8, 11 and comes
% back at (0, 2.5): length 11, s = -1 (one row lower; offsets measured
% upwards). The closing segment of its curve runs up the first side from 2.5
% to 3.5; continued along the walk, the chart brings x_0 back at x_0' =
% (11, 4.5), where the walk round T (keeping T on its right) starts.
% First visits on that walk, with R the quarter-turn x -> (-y, x), R_i = R for
% valence 3 and R^{-1} for valence 5, Q_i = R_6 ... R_{i+1}:
%   p_1 = (11,5) v3  (I-R)p = (16,-6)  Q_1 = R^3  -> ( -6,-16)
%   p_2 = ( 8,5) v3  (I-R)p = (13,-3)  Q_2 = R^2  -> (-13,  3)
%   p_3 = ( 6,5) v3  (I-R)p = (11,-1)  Q_3 = R    -> (  1, 11)
%   p_4 = ( 5,5) v5  (I-R^-1)p = (0,10) Q_4 = R^2 -> (  0,-10)
%   p_5 = ( 4,5) v3  (I-R)p = ( 9, 1)  Q_5 = R    -> ( -1,  9)
%   p_6 = ( 3,5) v3  (I-R)p = ( 8, 2)  Q_6 = I    -> (  8,  2)
%   B_D = sum = (-11, -1) = x_0 - x_0'.
% D lies on the left of the walk, so B_D = T_gamma = (-alpha, s): minus the
% length 11 along the walk, the shift -1 across it, s = <e_2, B_D>
% (Proposition 2(a)). The other order of visit gives
% (1,1), which is not a translation of x_0: the order matters.
% gamma' from (0, 1.5) comes back at (0, 0.5): s = -1, length 11. The curves of
% gamma and gamma' bound an annulus without singular vertex (shaded), and the
% shifts agree (Proposition 2(c)).
% Parity on a sphere (Proposition 5): 11 = -1 mod 2.
% check.py builds this surface and checks these numbers by walking the strips
% of the quad mesh: the valences above, every other vertex of valence 4,
% s = -1 and length 11 from o = 1.5, 2.5, 3.5, and the sum B_D.
\begin{tikzpicture}[x=0.55cm, y=0.55cm, font=\footnotesize,
    side/.style={draw=black, line width=0.5pt},
    quadgrid/.style={draw=black!18, line width=0.25pt, step={(1,1)}},
    surface/.style={fill=black!3},
    node dot/.style={circle, fill=black, inner sep=0pt, minimum size=2.6pt},
    turn/.style={line width=1pt, -{Stealth[length=5pt]}},
    closing/.style={line width=0.7pt, densely dotted},
    cut/.style={draw=black, line width=1.1pt, line join=round},
    brace/.style={decorate, decoration={brace, amplitude=3pt}, line width=0.5pt},
    burgers/.style={draw=measured, line width=1.1pt, -{Stealth[length=6pt]}},
    val3/.pic={\fill[valthree, draw=black, line width=0.3pt]
      (90:3.4pt) -- (210:3.4pt) -- (330:3.4pt) -- cycle;},
    val5/.pic={\fill[valfive, draw=black, line width=0.3pt]
      (90:3.4pt) -- (162:3.4pt) -- (234:3.4pt) -- (306:3.4pt) -- (18:3.4pt) -- cycle;}]
  \definecolor{loopspiral}{HTML}{D55E00}
  \definecolor{annulus}{HTML}{CFE0F3}
  \definecolor{valthree}{HTML}{E8890C}
  \definecolor{valfive}{HTML}{7C3AAD}
  \definecolor{measured}{HTML}{14934A}

  % The disc Omega: band, block, top.
  \fill[surface] (0,-0.6) rectangle (11,4) (3,4) rectangle (11,5)
    (0,4) rectangle (3,6) (0,6) rectangle (2,7);
  % The annulus between the curves of gamma' and gamma.
  \fill[annulus] (0,0.5) -- (0.15,0.5) -- (0.15,1.5) -- (11,1.5) -- (11,3.5)
    -- (0.15,3.5) -- (0.15,2.5) -- (0,2.5) -- cycle;
  \draw[quadgrid] (0,-0.6) grid (11,4);
  \draw[quadgrid] (3,4) grid (11,5);
  \draw[quadgrid] (0,4) grid (3,6);
  \draw[quadgrid] (0,6) grid (2,7);
  % Traces. The cut T: the left side up to c_1 and the border of the top; its
  % other side: the right side and the top of the block.
  \draw[side] (0,-0.6) -- (0,3.5) (11,-0.6) -- (11,4.5);
  \draw[cut] (0,3.5) -- (0,7) -- (2,7) -- (2,6) -- (3,6) -- (3,5) -- (11,5) -- (11,4.5);
  \draw[side] (0,4) -- (2,4) (2,6) -- (2,-0.6) (3,5) -- (2,5) (3,5) -- (3,-0.6) (2,6) -- (0,6);
  \foreach \x in {4,5,6,8} \draw[side] (\x,5) -- (\x,-0.6);
  \foreach \p in {(2,4), (2,5), (0,6), (9,5)} \node[node dot] at \p {};
  % The turns, their closing segments and where they come back.
  \foreach \o/\name in {3.5/{\gamma}, 1.5/{\gamma'}} {
    \draw[turn, loopspiral] (0,\o) -- (11.1,\o);
    \draw[closing, loopspiral] (0.15,\o-1) -- (0.15,\o);
    \draw[loopspiral, line width=0.5pt, fill=white] (0,\o-1) circle (1.8pt);
    \draw[loopspiral, line width=0.5pt, fill=white] (0,\o) circle (1.8pt);
    \node[above, loopspiral] at (9.6,\o) {$\name$};
    \node[right, loopspiral] at (11.2,\o) {$s = -1$};
  }
  \draw[loopspiral, line width=0.5pt, fill=white] (11,4.5) circle (1.8pt);
  \node[left=1pt] at (0,3.5) {$x_0$};
  \node[right=1pt] at (11,4.5) {$x_0'$};
  % The singular vertices: first visits p_i on the top of the block, and
  % the other copies on the border of the top.
  \pic at (3,5) {val3};
  \foreach \p in {(0,4), (11,5), (0,7), (8,5), (2,7), (6,5), (3,6), (4,5)} \pic at \p {val3};
  \foreach \p in {(2,6), (5,5)} \pic at \p {val5};
  \begin{scope}[font=\scriptsize, inner sep=1pt]
    \foreach \i/\x in {1/11, 2/8, 3/6, 4/5, 5/4} \node[above=2pt] at (\x,5) {$p_{\i}$};
    \node[above right=1pt] at (3,5) {$p_6$};
    \node[left=2pt] at (0,4) {$c_1$};
    \node[left=2pt] at (0,7) {$c_2$};
    \node[above=2pt] at (2,7) {$c_3$};
    \node[above left=0pt] at (2,6) {$c_4$};
    \node[right=2pt] at (3,6) {$c_5$};
  \end{scope}
  \draw[brace] (-1.5,3.5) -- (-1.5,7) node[midway, left=3pt] {$D$};
  \node[above] at (1,7) {$T$};
  % The Burgers vector of D, from x_0' to x_0, and its two components.
  \draw[burgers] (11,-1.6) -- (0,-2.6);
  \draw[measured, dashed, line width=0.5pt] (11,-1.6) -- (0,-1.6) -- (0,-2.6);
  \node[measured, above] at (5.5,-1.6) {along: $\alpha_\gamma = 11$};
  \node[measured, left] at (0,-2.1) {across: $s_\gamma = -1$};
  \node[measured, below] at (6.5,-2.3) {$B_D = \sum_i Q_i (I - R_i)\, p_i = (-11, -1)$};
\end{tikzpicture}

\caption{The shift read on the singularities
(Proposition~\ref{prop:burgers}). A cap $D$ with six singular vertices
(valences $3, 3, 3, 5, 3, 3$, total index $4$), cut along a tree $T$ and
developed; each $c_i$ is drawn at its first visit $p_i$ on the walk round $T$.
The turn $\gamma$, of length $11$, comes back one row lower, and $D$ lies on
its left: $B_D = \sum_i Q_i (I - R_i) p_i = (-11, -1) = T_\gamma(\ell)$,
minus the length along the walk and the shift $-1$ across it. The turn $\gamma'$ bounds with $\gamma$ an
annulus without singular vertex (shaded) and has the same shift (part~(c));
$11 \equiv -1 \pmod 2$ (Proposition~\ref{prop:parity}). In all figures
offsets are measured upwards.}
\label{fig:holonomy}
\end{figure}

The order matters: on Figure~\ref{fig:holonomy}, walking round $T$ the other
way gives $(1, 1)$ in place of $B_D = (-11, -1)$.

\begin{proof}
(a) Orient each edge $f$ of $T$ away from $x_0$; its two sides on the border of
$\Omega$ are glued by a motion $G_f$, from the side the walk round $T$ follows
first to the other. Let $f$ end at $v$, from which $f_1, \dots, f_r$ leave in
the order of the walk. A small arc turning clockwise round $v$ from the first
side of $f$ to the second crosses $f_1, \dots, f_r$ and turns through the
whole angle at $v$, so $G_f = G_{f_r} \circ \dots \circ G_{f_1} \circ
\rho_v$, with $\rho_v(x) = R^{4-k(v)}(x - p_v) + p_v$, $p_v$ the position at
the first visit ($\rho_v = I$ if $v$ is not singular). At $x_0$, with
$\bar\gamma$ run with $D$ on its right in place of the arc, $\hat
h(\bar\gamma) = G_{f_r} \circ \dots \circ G_{f_1}$. Unfolding the
recursion composes the $\rho_{c_i}$ in the order of first visit; the rotation
is $R^{\sum_i(4 - k(c_i))} = I$ by Lemma~\ref{lem:gauss}, expanding gives
$B_D = \hat h(\bar\gamma) = T_\gamma(\ell)^{-1}$, and with $D$ on its left
$\bar\gamma$ has the inverse holonomy, $B_D = T_\gamma(\ell)$. (b) With $x_0$ on $\bar\gamma_1$ at the start of $\kappa$, the loop
$\lambda = \bar\gamma_1\, \kappa\, \bar\gamma_2^{-1} \kappa^{-1}$ borders the
disc $A$ cut along $\kappa$. Developing along $\kappa$ reads $\hat
h(\bar\gamma_2)$ in the coordinates at $x_0$, where the walk of $\gamma_2$
runs along $e_1$; so $\hat h(\lambda)$ is the translation by $\hat
h(\bar\gamma_1) - \hat h(\bar\gamma_2)$, whose component along $e_2$ is
$\shift_{\gamma_2}(\ell) - \shift_{\gamma_1}(\ell)$, and part~(a) applies to the cut disc,
whose singular vertices have total index zero. (c) The sum is empty.
\end{proof}

In the T-meshes of openclare, whose traces start at singular vertices, an
annulus without singular vertex holds no node unless a trace winding round
enters it through a closing segment, so two turns as in (c) usually lie in
one corridor.

The simplest configuration between two parallel turns is a pair of valences
three and five, which lets the number of edges round a limb change. A pair of
opposite disclinations of rotation $\Omega$ is a dislocation with a Burgers
vector of $2\sin(\Omega/2)$ times their
separation~\cite[Eq.~(5)]{kleman2008disclinations}. We read it
on the shifts of two turns (Figures~\ref{fig:holonomy}
and~\ref{fig:dislocation}).

\begin{proposition}[A dislocation between two turns]\label{prop:dislocation}
In the setting of Proposition~\ref{prop:burgers}(b), let the annulus $A$ hold
exactly two singular vertices, $c_3$ of valence three and $c_5$ of valence
five, and let $d = p_5 - p_3 = d_1 e_1 + d_2 e_2$ be their separation,
developed inside the cut disc. Then:
\begin{enumerate}
\item[(i)] $\shift_{\gamma_1}(\ell) - \shift_{\gamma_2}(\ell) = \varepsilon_1
  d_1 + \varepsilon_2 d_2$ with signs $\varepsilon_1, \varepsilon_2 \in \{+1,
  -1\}$ that depend on the configuration but not on $\ell$, and $|B_A| =
  \sqrt2\,|d|$. The two turns can close together only if $c_5$ lies, as seen
  from $c_3$, on one particular diagonal of the field.
\item[(ii)] If $c_3$ and $c_5$ lie on one straight line of the field inside
  $A$, of length $\delta$ (in a T-mesh: they are joined by a chain of arcs
  along one trace, and $\delta$ is the sum of their lengths), and some path
  $\kappa$ as in Proposition~\ref{prop:burgers}(b) does not cross that line,
  then
  $|\shift_{\gamma_1}(\ell) - \shift_{\gamma_2}(\ell)| = \delta$. The two turns
  can both be rings only if $\delta = 0$, that is if the two singularities
  merge; as long as the bounds of separation keep $\delta \ge 1$, one of them at
  least is a helix.
\item[(iii)] For whole lengths, $\shift_{\gamma_1}(\ell) -
  \shift_{\gamma_2}(\ell) \equiv d_1 + d_2 \pmod 2$: a pair at an odd number of
  edges from one another, counted along the two directions of the field, never
  lies between two rings.
\end{enumerate}
\end{proposition}

\begin{proof}
With $R_3 = R$ and $R_5 = R^{-1}$, the two orders give $B_A = R^{-1}(I -
R)p_3 + (I - R^{-1})p_5 = (I - R^{-1})d$ or $B_A = R(I - R^{-1})p_5 + (I -
R)p_3 = (I - R)(-d)$. With $R(a, b) = (-b, a)$, $(I - R^{-1})(d_1, d_2) = (d_1 -
d_2,\ d_1 + d_2)$ and $(I - R)(-d_1, -d_2) = (-d_1 - d_2,\ d_1 - d_2)$;
reversing the orientation of the border inverts the rotations and gives the
two other sign patterns (checked symbolically). In all four both components
are $\pm d_1 \pm d_2$ and $|B_A|^2 = 2|d|^2$, which is (i). For (ii), take
$\kappa$ not crossing the segment joining $c_3$ to $c_5$ (crossing it would
turn the frame by a quarter-turn round $c_3$) and a tree containing that
segment, and develop $c_5$ from
$c_3$ along it: $d$ is $\delta$ times one of $\pm e_1$, $\pm e_2$. For (iii), $d
\in \Z^2$ and $\pm d_1 \pm d_2 \equiv d_1 + d_2$.
\end{proof}

\begin{figure}[t]
\centering
% A 3-5 dislocation between two parallel turns (Propositions 3 and 6).
% \input inside a figure; needs \usepackage{tikz} and
% \usetikzlibrary{arrows.meta,decorations.pathreplacing}.
%
% Each panel is a band of a T-mesh on a cylinder, cut open along a trace of
% the vertex c_5 of valence five: the right side is glued to the left side, the
% lower part (below c_5) at the same height, the upper part two units higher.
% Every node is at a whole position and the faint grid is the quad mesh at
% these lengths. c_5 is drawn three times (its angle 5 x 90 degrees is split
% between the three copies: 90 at the bottom left, 180 at the right, 180 at
% the left), c_3 once (270 degrees). The empty corner next to c_3 is not part of
% the surface: its two thick sides are glued to each other by a quarter-turn
% about c_3. The turn gamma_1 (below the pair) and gamma_2 (above it) bound an
% annulus that holds exactly c_3 and c_5 (Proposition 2(b)).
%
% (a) Straight chain (Proposition 3(ii)). c_3 = (2,3), c_5 = (0,3) = (10,3) =
% (2,5); the gluings are (10,y) ~ (0,y) for y < 3, (10,y) ~ (2,y+2) for y > 3,
% and the quarter-turn about c_3 taking [c_3, (0,3)] onto [c_3, (2,5)]. c_3 and
% c_5 are joined along the trace y = 3 by the chain of arcs b = [c_3, (1,3)]
% and c = [(1,3), c_5]; a trace from below (its vertex is out of the picture)
% stops at (1,3). delta = l_b + l_c = 1 + 1 = 2.
%   gamma_1 at o = 1.5: crosses x = 1, 2, 10; comes back at 1.5: s = 0,
%   length 10 (a ring).
%   gamma_2 at o = 5.5 on the side x = 2: crosses x = 10 at height 5.5, which
%   is height 7.5 on the side x = 2: s = 2, length 10 - 2 = 8 (a helix).
%   |s_2 - s_1| = 2 = delta (Proposition 3(ii)); the lengths differ by 2 too:
%   the Burgers vector of the pair is (2, 2), of norm sqrt(2) |d| with d = (2, 0)
%   from c_3 to c_5 (Proposition 3(i)); s_2 - s_1 = 2 = d_1 + d_2 mod 2 (iii).
%   The separatrix y = 5 of c_5 goes once round and stops on the trace x = 10
%   two units above c_5 (Proposition 8 gives s_2 = 2 again).
%   Both turns can be rings only if delta = 0: c_3 and c_5 merge.
%   Parity on a sphere (Proposition 5): 10 = 0 and 8 = 2 mod 2.
%
% (b) Staircase (Proposition 6). c_3 = (1,4), c_5 = (0,3) = (10,3) = (0,5);
% gluings (10,y) ~ (0,y) for y < 3, (10,y) ~ (0,y+2) for y > 3, and the
% quarter-turn about c_3 taking [c_3, (1,3)] onto [c_3, (0,4)] and
% [(1,3), (0,3)] onto [(0,4), (0,5)]. c_3 and c_5 are on no common trace; the
% staircase from c_3 to c_5 is b = [c_3, (1,3)], down, then c = [(1,3), c_5],
% left: d = p_5 - p_3 = (-1, -1), and with epsilon = (-1, -1) each step has
% epsilon . v = 1 (b: v = (0,-1); c: v = (-1,0)).
%   gamma_1 at o = 1.5: s = 0, length 10. gamma_2 at o = 5.5: crosses x = 10
%   at height 5.5 = height 7.5 on the side x = 0: s = 2, length 10.
%   s_2 - s_1 = 2 = l_b + l_c = epsilon . d (Proposition 6); here the lengths
%   are equal, the Burgers vector of the pair is (0, 2). Both are rings only
%   if l_b = l_c = 0. Parity: 10 = 0 = 2 mod 2.
% check.py builds both bands and checks these numbers by walking the strips
% of the quad mesh: valences 5 and 3 at the copies listed above, every other
% vertex of valence 4, and the turns above.
\begin{tikzpicture}[x=0.5cm, y=0.5cm, font=\footnotesize,
    side/.style={draw=black, line width=0.5pt},
    quadgrid/.style={draw=black!18, line width=0.25pt, step={(1,1)}},
    surface/.style={fill=black!3},
    node dot/.style={circle, fill=black, inner sep=0pt, minimum size=2.6pt},
    arc mark/.style={line width=1.8pt, draw=black},
    turn/.style={line width=1pt, -{Stealth[length=5pt]}},
    brace/.style={decorate, decoration={brace, amplitude=3pt}, line width=0.5pt},
    val3/.pic={\fill[valthree, draw=black, line width=0.3pt]
      (90:3.4pt) -- (210:3.4pt) -- (330:3.4pt) -- cycle;},
    val5/.pic={\fill[valfive, draw=black, line width=0.3pt]
      (90:3.4pt) -- (162:3.4pt) -- (234:3.4pt) -- (306:3.4pt) -- (18:3.4pt) -- cycle;}]
  \definecolor{loopspiral}{HTML}{D55E00}
  \definecolor{loopclosed}{HTML}{0B3C8C}
  \definecolor{valthree}{HTML}{E8890C}
  \definecolor{valfive}{HTML}{7C3AAD}

  % ---------------------------------------------------------------- (a)
  \begin{scope}
    % The surface: the lower band [0,10]x[0,3] and the upper part [2,10]x[3,8].
    \fill[surface] (0,-0.6) rectangle (10,3) (2,3) rectangle (10,8.6);
    \draw[quadgrid] (0,-0.6) grid (10,3);
    \draw[quadgrid] (2,3) grid (10,8.6);
    % Traces: the cut (a trace of c_5, both copies), the trace from c_3 round
    % to c_5, c_3 downwards, a trace from below, the winding separatrix.
    \draw[side] (0,-0.6) -- (0,3) (10,-0.6) -- (10,8.6) (2,5) -- (2,8.6);
    \draw[side] (2,3) -- (10,3) (2,-0.6) -- (2,3) (1,-0.6) -- (1,3) (2,5) -- (10,5);
    % The chain from c_3 to c_5, and its glued copy.
    \draw[arc mark] (0,3) -- (2,3) (2,3) -- (2,5);
    \node[below] at (1.5,3) {$b$};
    \node[below] at (0.5,3) {$c$};
    \foreach \p in {(1,3), (2,4), (10,5), (2,7)} \node[node dot] at \p {};
    % The turns.
    \draw[turn, loopclosed] (0,1.5) -- (10.1,1.5);
    \draw[loopclosed, line width=0.5pt, fill=white] (0,1.5) circle (1.8pt);
    \node[above, loopclosed] at (6,1.5) {$\gamma_1$};
    \draw[turn, loopspiral] (2,5.5) -- (10.1,5.5);
    \draw[loopspiral, line width=0.5pt, fill=white] (2,5.5) circle (1.8pt);
    \draw[loopspiral, dashed, line width=0.8pt, -{Stealth[length=4pt]}] (2,7.5) -- (4,7.5);
    \draw[loopspiral, line width=0.5pt, fill=white] (2,7.5) circle (1.8pt);
    \node[above, loopspiral] at (6,5.5) {$\gamma_2$};
    \draw[brace] (1.7,5.5) -- (1.7,7.5) node[midway, left=3pt] {$2$};
    % The singular vertices.
    \pic at (2,3) {val3};
    \foreach \p in {(0,3), (10,3), (2,5)} \pic at \p {val5};
    \node[above right=-1pt] at (2,3) {$c_3$};
    \node[left=1pt] at (0,3) {$c_5$};
    % Right: shifts and lengths.
    \node[right, align=left, loopspiral] at (10.3,5.5) {$s = 2$\\ 8 quads};
    \node[right, align=left, loopclosed] at (10.3,1.5) {$s = 0$\\ 10 quads};
    \node[anchor=north, align=center] at (6,-1.1)
      {$s_{\gamma_2} - s_{\gamma_1} = \ell_b + \ell_c = \delta = 2$};
    \node[anchor=north] at (6,-2.4) {\small (a) straight chain};
  \end{scope}

  % ---------------------------------------------------------------- (b)
  \begin{scope}[shift={(16.8,0)}]
    % The surface: [0,10]x[0,3], [1,10]x[3,4], [0,10]x[4,8].
    \fill[surface] (0,-0.6) rectangle (10,3) (1,3) rectangle (10,4) (0,4) rectangle (10,8.6);
    \draw[quadgrid] (0,-0.6) grid (10,3);
    \draw[quadgrid] (1,3) grid (10,4);
    \draw[quadgrid] (0,4) grid (10,8.6);
    % Traces: the cut (a trace of c_5), the loop of c_5 at y = 3, the three
    % separatrices of c_3, the winding separatrix of c_5.
    \draw[side] (0,-0.6) -- (0,3) (10,-0.6) -- (10,8.6) (0,5) -- (0,8.6);
    \draw[side] (0,3) -- (10,3) (1,4) -- (10,4) (1,4) -- (1,5) (0,5) -- (10,5);
    % The staircase from c_3 to c_5, and its glued copy.
    \draw[arc mark] (1,4) -- (1,3) -- (0,3) (1,4) -- (0,4) -- (0,5);
    \node[right] at (1,3.5) {$b$};
    \node[below] at (0.5,3) {$c$};
    \foreach \p in {(1,3), (0,4), (10,4), (0,6), (1,5), (10,5), (0,7)}
      \node[node dot] at \p {};
    % The turns.
    \draw[turn, loopclosed] (0,1.5) -- (10.1,1.5);
    \draw[loopclosed, line width=0.5pt, fill=white] (0,1.5) circle (1.8pt);
    \node[above, loopclosed] at (6,1.5) {$\gamma_1$};
    \draw[turn, loopspiral] (0,5.5) -- (10.1,5.5);
    \draw[loopspiral, line width=0.5pt, fill=white] (0,5.5) circle (1.8pt);
    \draw[loopspiral, dashed, line width=0.8pt, -{Stealth[length=4pt]}] (0,7.5) -- (2,7.5);
    \draw[loopspiral, line width=0.5pt, fill=white] (0,7.5) circle (1.8pt);
    \node[above, loopspiral] at (6,5.5) {$\gamma_2$};
    \draw[brace] (-0.3,5.5) -- (-0.3,7.5) node[midway, left=3pt] {$2$};
    % The singular vertices.
    \pic at (1,4) {val3};
    \foreach \p in {(0,3), (10,3), (0,5)} \pic at \p {val5};
    \node[above right=-1pt] at (1,4) {$c_3$};
    \node[left=1pt] at (0,3) {$c_5$};
    % Right: shifts and lengths.
    \node[right, align=left, loopspiral] at (10.3,5.5) {$s = 2$\\ 10 quads};
    \node[right, align=left, loopclosed] at (10.3,1.5) {$s = 0$\\ 10 quads};
    \node[anchor=north, align=center] at (6,-1.1)
      {$s_{\gamma_2} - s_{\gamma_1} = \ell_b + \ell_c = 2$};
    \node[anchor=north] at (6,-2.4) {\small (b) staircase};
  \end{scope}
\end{tikzpicture}

\caption{A 3--5 dislocation between a ring $\gamma_1$ and a helix $\gamma_2$
of a band, cut open along a trace of the vertex of valence five, which is
drawn three times. (a)~The pair is joined along one trace by two arcs, $\delta
= \ell_b + \ell_c = 2$: $\gamma_2$ comes back $2$ rows higher, with $8$
quads per turn, and the Burgers vector of the pair is $(2, 2)$
(Proposition~\ref{prop:dislocation}). (b)~The pair is joined by a staircase of
two arcs, $d = \varepsilon = (-1, -1)$: $\gamma_2$ comes back $2$ rows higher,
with $10$ quads (Proposition~\ref{prop:staircase}). In both, the shifts differ
by $\ell_b + \ell_c = 2$, and both turns are rings only if the two vertices
merge.}
\label{fig:dislocation}
\end{figure}

\begin{proposition}[Checkerboard rule]\label{prop:checkerboard}
In the setting of Proposition~\ref{prop:burgers}(b), with whole lengths,
colour the points $p \in \Z^2$ of the cut disc black or white by the parity of
$p_1 + p_2$ (either colour may be called white). Then $\shift_{\gamma_1}(\ell)
- \shift_{\gamma_2}(\ell)$ has the parity of the number of singular vertices
of odd valence inside $A$ that develop onto white points. In particular, two
rings bound an annulus only if that number is even.
\end{proposition}

\begin{proof}
In the sum of Proposition~\ref{prop:burgers}, a vertex of even valence has
$R_i = \pm I$ and $(I - R_i)p_i \in \{0, 2p_i\}$; for odd valence both
components of $(I - R^{\pm1})(a, b) = (a \pm b,\ b \mp a)$ have the parity of
$a + b$. Quarter-turns keep the property that both components have a given
parity, and sums add the parities, which gives the count. A quarter-turn keeps
colours; a whole translation with $t_1 + t_2$ odd swaps them, which keeps the
parity of the count because the vertices of odd valence inside $A$ are even in
number, their odd indices summing to zero with the even ones
(Lemma~\ref{lem:gauss}). Developing a vertex along another path inside the cut
disc rotates it about another singular vertex, which keeps its colour.
\end{proof}

The rule also holds on $13{,}000$ random configurations checked with SymPy.

\paragraph{Parity on a sphere.}
Each turn of a closed sphere carries its own parity condition, which needs
neither a simple curve nor the holonomy.

\begin{proposition}[Parity of a turn on a sphere]\label{prop:parity}
Let the surface be closed of genus zero and the T-mesh admit lengths that are
positive and balance every patch (its real lengths do). For every turn
$\gamma$ and all whole lengths $\ell$ balancing every patch, of any sign,
\[
  \alpha_\gamma(\ell) \equiv \shift_\gamma(\ell) \pmod 2 .
\]
\end{proposition}

\begin{proof}
First let $\ell$ be whole, balanced and positive: the patches form a quad mesh
of the sphere, whose graph is bipartite, since $H_1(S^2; \Z/2) = 0$ makes
every cycle a sum of faces of four edges. In a 2-colouring the colour of the
point $(i, j)$ of a patch $P$ is $i + j + \kappa_P \bmod 2$. For $M(x) = R^k
x + t$, $t \in \Z^2$, the coordinate sum of $M(x)$ has the parity of that of
$x$ plus $\pi(M) = t_1 + t_2$. Let $P_0, \dots, P_N = P_0$ be the patches
crossed by $\gamma$ and $D_k$ the motion developing $P_k$ into $P_0$, with
$D_N = T_\gamma(\ell)$. A node of the arc from $P_k$ to $P_{k+1}$, of
coordinates $y$, $y'$ with $D_k(y) = D_{k+1}(y')$, has the colour $(D_k y)_1 +
(D_k y)_2 + \pi(D_k) + \kappa_{P_k}$, and likewise from $P_{k+1}$; so
$\kappa_{P_k} + \pi(D_k)$ is constant, and at $k = N$, $\pi(T_\gamma(\ell))
\equiv 0$. In the frame of the first side, that of $P_0$ up to a quarter-turn
and signs, $T_\gamma(\ell) = (\alpha_\gamma(\ell),
\pm\shift_\gamma(\ell))$. For general whole balanced $\ell$: the balanced
lengths form a rational subspace with positive real points, so there are
positive whole balanced lengths $p$, and the linear form $\alpha_\gamma +
\shift_\gamma$ takes at $\ell = (\ell + Np) - Np$, $N$ large, the difference
of two even values.
\end{proof}

The proof uses only that the quad graph is bipartite; it fails in higher genus
(a $3 \times 3$ grid on a torus has rings of three quads) and with borders
(Theorem~\ref{thm:parity-any}). On a closed sphere it gives at once:

\begin{corollary}\label{cor:parity}
On a closed surface of genus zero:
\begin{enumerate}
\item[(a)] every ring of quads has an even number of quads;
\item[(b)] a helix of odd pitch winds round an odd number of quads per turn;
\item[(c)] closing a helix of odd pitch changes the length of its walk by an
  odd number of edges: a repair that keeps the length of a turn cannot close it
  if its pitch is odd, and changing the pitch by an odd amount always changes
  the length.
\end{enumerate}
\end{corollary}

\begin{proof}
(a) A ring is the strip of a walk with $\shift =
0$~\cite[Prop.~2]{gardent2026loops}, so $\alpha$ is even. (b) and (c) are the
proposition read before and after the change.
\end{proof}

Part (a) is the classical bipartiteness of quadrangulations of the
sphere~\cite[Lemma~3.1]{kaiser2015colouring}. For a turn bounding a disc, the
argument of the checkerboard rule gives both components of $B_D =
(\alpha_\gamma, \shift_\gamma)$ the parity of the number of vertices of odd
valence on white points, which is the proposition for that turn.

\begin{remark}[Borders; a second proof]\label{rem:borders}
The proposition and part (a) of the corollary need the surface to be closed:
an open cylinder of three quads round is a ring of three quads. A second proof,
from the formalization, needs neither the bipartite graph nor positive
lengths: every loop of $\Sigma^\circ$ is a product of loops round the
singular vertices (Section~\ref{sec:fox}), quarter-turns about whole points
that keep the parity of $p_1 + p_2$, so the translation of every loop has an
even coordinate sum.
\end{remark}

\paragraph{How many independent rings.}
Rings that do not cross are disjoint simple closed curves of $\Sigma^\circ$;
on a closed sphere with $n \ge 3$ singular vertices at most $n - 3$ are
pairwise non-homotopic, as for the cylinders of a Jenkins--Strebel
differential~\cite{strebel1984quadratic}, and homotopic ones share their shift,
so their closing conditions have rank at most $n - 3$ ($3g - 3 + n + b$ plus at
most $b$ classes of collars with genus $g$ and $b$ borders).

\paragraph{Certificates.}
Closing a set of turns within bounds is a question of linear feasibility.

\begin{remark}[Sign-definite turns; Farkas certificates]\label{rem:definite}
If the shift of a turn equals, modulo the balances, a form $\sum_a c_a \ell_a$
with every $c_a \ge 0$, then for all balanced $\ell \ge 0$ the shift
$\shift_\gamma(\ell)$ is never negative, and it is zero only if every arc with
$c_a > 0$ has length zero. If one of those arcs is held open by a bound of
separation, the turn is a helix under every quantization meeting the bounds;
it should never be held, and one linear program per turn finds such turns.
More generally, by Farkas' lemma (with $x = b + z$, $z \ge 0$), the system
$B\ell = 0$, $\shift_\gamma(\ell) = 0$ for $\gamma \in H$, $\ell \ge b$ has no
solution exactly when some combination $\sum_{\gamma \in H} u_\gamma
\shift_\gamma$ equals, modulo the balances, a form $\sum_a c_a \ell_a$ with $c
\ge 0$ and $\sum_a c_a b_a > 0$; and an arc $a$ has length zero in every
solution with $\ell \ge 0$ exactly when some such combination has $c_a > 0$.
The data being integral, bounds included, the multipliers of the shifts and of
the balances can be taken whole.
\end{remark}

\begin{remark}[Whole solutions exist at some size]\label{rem:exist}
The balance and closing constraints are homogeneous and the bounds of
separation are lower bounds $\ell \ge b$ with $b \ge 0$. If they have a real
solution they have a rational one, a vertex of a rational polyhedron, and
multiplying it by the common denominator $N$ gives whole lengths $N\ell \ge Nb
\ge b$ that satisfy them all. Closing in whole numbers is thus never
impossible in itself; what can fail is closing \emph{at a given size} (number
of quads) or within a given deviation from the targets.
\end{remark}

\paragraph{On the male body.}
The statements of this section hold wherever they apply (measured). The
crossing maps give the exported shift of all $1{,}299$ turns. On the
$1{,}436$ sides of the $718$ simple ones, of index $4$ and $4$ to $257$
singular vertices, the holonomy equals the sum of
Proposition~\ref{prop:burgers}(a) in both components, while the order of
last visit fails on $150$ of the $158$ discs of way~2; the checkerboard rule
holds there and on the $30{,}148$ annuli between disjoint simple turns, of
index $0$; the $294$ pairs with no singular vertex between them have equal
shifts; on the ten annuli holding one 3--5 pair,
Proposition~\ref{prop:dislocation}(i) holds, the two never on one trace.
Proposition~\ref{prop:parity} holds on all turns and on the $1{,}838$ ways
round of the strips of the eight meshes. Of the $101$ turns crossing themselves
once, $97$ split as $(3, 3, 2)$ and $4$, crossing their closing segment, as
$(2, 2, 4)$; $480$ cross themselves up to $107$ times, and only $356$ of the
$765$ held turns are simple: Section~\ref{sec:fox} handles any curve.

\section{Conflicts between turns}\label{sec:conflicts}

The held turns of the body close together if lengths are only asked to be
non-negative, not within $P_{\mathrm{eng}}$. Table~\ref{tab:iis} lists
\emph{irreducible} infeasible sets, which cannot be closed within the bounds
while every proper subset can; the Farkas multipliers were checked in exact
arithmetic, and for each turn left out, whole lengths closing the others were
found by HiGHS and checked in exact integer arithmetic (measured).

\begin{table}[t]
\centering
\small
\begin{tabular}{llllll}
\toprule
ways & size & turns & multipliers & bounded arcs & exact check \\
\midrule
0, 1, 2, 4 & 1 & 42 (way 0) & $-1$ & 1937 & way 0 \\
1, 3, 5, 7 & 2 & 54, 61 (way 1) & $1, -1$ & 1587 & --- \\
1, 3 & 2 & 45, 49 (way 1) & $1, -1$ & 1069 & --- \\
1, 3, 7 & 2 & 26, 37 (way 1) & $-1, 1$ & 619 & --- \\
5, 7 & 2 & 58, 94 (way 5) & $1, -1$ & 858 & --- \\
5 & 2 & 49, 71 & $-1, 1$ & 619, 625 & --- \\
6 & 2 & 57, 58 & $1, 1$ & 1241 & --- \\
7 & 2 & 8, 15 & $1, 1$ & 191 & --- \\
1 & 3 & 36, 37, 39 & $1, 1, -1$ & 619 & yes \\
1 & 4 & 40, 56, 59, 61 & $1, 1, -1, -1$ & 1587 & yes \\
1 & 4 & 66, 71, 75, 76 & $-\tfrac76, -\tfrac56, -\tfrac13, -\tfrac12$ & 1967, 1980, 1998 & yes \\
3 & 4 & 73, 78, 82, 83 & $-\tfrac74, -\tfrac54, -\tfrac12, -\tfrac34$ & 1980, 1998 & yes \\
5 & 7 & 59, 91, 96, 98, 99, 157, 158 & $2, 4, 1, -4, -1, 1, -1$ & 1036 & yes \\
7 & 7 & 112--116, 145, 147 & $8, \tfrac12, 1, \tfrac12, \tfrac12, -\tfrac12, \tfrac12$ & 1157 & yes \\
5 & 15 & 0, 2--9, 11, 14, 16--18, 64 & whole, up to 23 & 17 & 14 of 15 \\
\bottomrule
\end{tabular}
\caption{Irreducible infeasible sets of held turns on the male body, bounds of
separation kept (measured). Turn numbers are those of the first way listed;
``bounded arcs'': arcs held at length at least one in the support of the
non-negative form. Rows without an exact check come from the floating-point
linear program; their turns are each closable alone, so a pair is
irreducible. In the last row, the whole lengths closing the set without turn 3
were not found within the time limit.}
\label{tab:iis}
\end{table}

In four ways the turn round the left wrist cannot even be closed alone: its
shift is $-(\ell_{1921} + \ell_{1937})$ modulo the balances, and arc $1937$,
$0.08$ edges long, is held at length at least one (exact for way 0; whether
every valid quantization meets this bound we did not check). One may hope that every other conflict contains two parallel
turns whose annulus forbids it. The smallest counterexample has three turns
near the right eye, every two of which close together:
$\shift_{36} + \shift_{37} - \shift_{39} = \ell_{604} + \ell_{619} +
\ell_{620} + \ell_{637}$ modulo the balances, with $\ell_{619} \ge 1$. Every
set found is infeasible in real numbers already, so none is an obstruction of
parity. The conflicts between two turns lie in clusters of singularities at the eyes,
hands and feet, the sum or difference of the shifts being a sum of two to
eight arc lengths, one bounded; yet in none of the $13$ we examined (the pairs of
Table~\ref{tab:iis} but $49$, $71$ of way~5) are the turns separated by a
3--5 pair on one trace, as in Proposition~\ref{prop:dislocation}(ii). What
separates them is a staircase.

\begin{proposition}[Staircases]\label{prop:staircase}
Let $\gamma_1$, $\gamma_2$ be turns and $c$, $c'$ singular vertices such that
\[
  \shift_{\gamma_1}(\ell) - \shift_{\gamma_2}(\ell) = \varepsilon_1 d_1 +
  \varepsilon_2 d_2
\]
for all balanced $\ell$, with $\varepsilon_1, \varepsilon_2 \in \{+1, -1\}$
and $d = p_{c'} - p_c = d_1 e_1 + d_2 e_2$ the separation of $c'$ from $c$
developed along a path $\pi$ from $c$ to $c'$
(Propositions~\ref{prop:dislocation}(i) and~\ref{prop:swap} give such
identities). Let $\sigma$ be a path from $c$ to $c'$ along arcs $a_1, \dots,
a_r$, through no other singular vertex, and homotopic to $\pi$ with its ends
fixed in $\Sigma$ without the other singular vertices. Developed along
$\sigma$, each arc $a_i$ runs along a direction $v_i \in \{\pm e_1, \pm
e_2\}$. If $\varepsilon_1 v_{i,1} + \varepsilon_2 v_{i,2} = 1$ for every $i$,
so that $\sigma$ is a \emph{staircase}, using two directions a quarter-turn
apart and climbing across the diagonal on which the two turns could close,
then
\[
  \shift_{\gamma_1}(\ell) - \shift_{\gamma_2}(\ell) = \sum_{i=1}^r \ell_{a_i}
\]
for all balanced $\ell$. Both turns are rings only if every arc of $\sigma$
has length zero. If this sum equals, modulo the balances, a form with
non-negative coefficients containing an arc held at length at least one, one
of the two turns is a helix under every quantization meeting the bounds.
\end{proposition}

\begin{proof}
Developing along $\sigma$ adds $\ell_{a_i} v_i$ arc after arc, the nodes passed
being T-junctions. As $\sigma$ and $\pi$ are homotopic with fixed ends away
from the other cone points, $d = \sum_i \ell_{a_i} v_i$, so $\varepsilon_1 d_1
+ \varepsilon_2 d_2 = \sum_i \ell_{a_i}$, first for positive balanced lengths,
then for all by linearity.
\end{proof}

The proof applies to any path on whose steps the form of the shift difference
is one; Proposition~\ref{prop:dislocation}(ii) is the straight case.
Exchanges give the other identities.

\begin{proposition}[Exchanging two vertices]\label{prop:swap}
Let $\gamma_1$, $\gamma_2$ be turns, and $\kappa$ a path from the start of
$\bar\gamma_1$ to that of $\bar\gamma_2$ that carries the walking direction of
$\gamma_2$ onto that of $\gamma_1$, as in Proposition~\ref{prop:burgers}(b);
let $\lambda = \bar\gamma_1\, \kappa\, \bar\gamma_2^{-1} \kappa^{-1}$.
Suppose that $\lambda$ is homotopic in $\Sigma^\circ$ to $\lambda_c
\lambda_{c'}^{-1}$, where $\lambda_c$ and $\lambda_{c'}$ are small loops run
with $c$, $c'$ on their right round two singular vertices of the same valence
$k$, reached along paths that develop them at $p_c$ and $p_{c'}$. Then
\[
  \shift_{\gamma_1}(\ell) - \shift_{\gamma_2}(\ell) = \bigl\langle e_2,\
  (I - R^{4-k})(p_{c'} - p_c) \bigr\rangle .
\]
For valence three or five this is $\pm d_1 \pm d_2$ with $d = p_c - p_{c'}$,
as for a dislocation, and Proposition~\ref{prop:staircase} applies.
\end{proposition}

\begin{proof}
As in Proposition~\ref{prop:burgers}(b), $\hat h(\lambda)$ is the translation by
$\hat h(\bar\gamma_1) - \hat h(\bar\gamma_2)$, whose component along $e_2$
is $\shift_{\gamma_2}(\ell) - \shift_{\gamma_1}(\ell)$. With $j = 4 - k$, $\hat
h(\lambda_c) \circ \hat h(\lambda_{c'})^{-1}(x) = R^j\bigl(R^{-j}(x - p_{c'}) +
p_{c'} - p_c\bigr) + p_c = x + (I - R^j)(p_c - p_{c'})$; for $j = \pm1$ both
components are $\pm d_1 \pm d_2$.
\end{proof}

The $13$ conflicts fall into six configurations (Table~\ref{tab:pairs},
measured). In A the annulus between two simple turns holds a 3--5 pair, and the
staircase joining them, $\ell_{611} + \ell_{624}$, $0.990$ long, is the
difference of the real shifts; modulo the balances it equals the certificate
$\ell_{611} + \ell_{616} + \ell_{619} + \ell_{635}$, which holds arc $619$
(exact);
at the engine's whole lengths the pair is one edge apart on one grid line,
Proposition~\ref{prop:dislocation}(ii) with $\delta = 1$ (its hypothesis on
$\kappa$ not checked). In B, C and D
the turns cross and their combination winds $+1$ and $-1$ round two vertices
of valence three, an exchange in homology (the homotopy of
Proposition~\ref{prop:swap} was not checked, the staircase identity was). In E
and F the combination winds once round vertices of total index four, like a
turn round a cap, and no staircase, nor sum of two, was found equal to it.

\begin{table}[t]
\centering
\small
\setlength{\tabcolsep}{4pt}
\newcolumntype{L}[1]{>{\raggedright\arraybackslash}p{#1}}
\begin{tabular}{@{}lL{4.1cm}L{3.9cm}L{4.4cm}L{2.0cm}@{}}
\toprule
& conflict & between the turns & staircase & mechanism \\
\midrule
A & $\shift_{37} - \shift_{26}$, ways 1, 3, 7; right eye; disjoint, simple
  & annulus: 78 (val.\ 3), 76 (val.\ 5)
  & 2 arcs, $0.990$; held arc 619 by the balances
  & dislocation \\
B & $\shift_{45} - \shift_{49}$, ways 1, 3; right hand; simple, crossing
  & lenses: 102, 103 (val.\ 3)
  & 4 arcs, $1.731$; held arc 1069 by the balances
  & exchange \\
C & $\shift_{58} - \shift_{94}$, ways 5, 7; right hand
  & 109, 118 (val.\ 3)
  & 3 arcs, $2.203$; held arc 858 by the balances
  & exchange \\
D & $\shift_{57} + \shift_{58}$, way 6; left foot; figure-eights
  & 155, 156 (val.\ 3)
  & 2 arcs, $1.452$, one of them held
  & exchange \\
E & $\shift_{54} - \shift_{61}$, ways 1, 3, 5, 7; left eye
  & 11 vertices of index 3, and 207 (val.\ 3)
  & none; a closed chain of 6 combs
  & combs \\
F & $\shift_{8} + \shift_{15}$, way 7; right foot
  & 14, 23, 24, and 15 (val.\ 3)
  & none; a closed chain of 4 combs round a cap
  & combs \\
\bottomrule
\end{tabular}
\caption{The conflicts between two turns on the male body (measured). Turn
numbers are those of the first way listed, signs those of the exported
shifts; ``between the turns'': the singular vertices, numbered as nodes,
round which the combination of the two curves winds; ``staircase'': one
between them whose length equals the combination modulo the balances, and
where the held arc of the certificate lies. A:
Propositions~\ref{prop:dislocation}(i) and~\ref{prop:staircase}; B--D:
Propositions~\ref{prop:swap} and~\ref{prop:staircase}; E, F:
Section~\ref{sec:combs}.}
\label{tab:pairs}
\end{table}

\paragraph{A separatrix that winds.}
The turn that cannot be closed alone is of another kind
(Figure~\ref{fig:separatrix}).

\begin{proposition}[A separatrix that winds]\label{prop:winding}
Let $\sigma$ and $\tau$ be separatrices of a singular vertex $v$ such that
$\sigma$, followed from $v$, stops on $\tau$ at a point $q$ at distance
$\delta(\ell)$ from $v$ along $\tau$, a sum of arc lengths. Let $\lambda$ be
the loop that runs along $\sigma$ from $v$ to $q$ and back along $\tau$ to $v$,
pushed off $v$ into the sector it turns through. If $\lambda$ is freely
homotopic in $\Sigma^\circ$ to the curve $\bar\gamma$ of a turn $\gamma$, the
walking direction of $\gamma$ going to that of $\sigma$, then
$\shift_\gamma(\ell) = \pm\delta(\ell)$ for all balanced $\ell$, the sign
depending on the side of $\sigma$ on which $\tau$ lies ($+$ in
Figure~\ref{fig:separatrix}(b), $-$ at the left wrist of the male body). The turn is a
ring only if $\delta = 0$, that is if $\sigma$ comes back onto its own vertex;
if $\delta$ equals, modulo the balances, a non-negative form containing an arc
held at length at least one, $\gamma$ is a helix under every quantization
meeting the bounds.
\end{proposition}

\begin{proof}
Develop $\lambda$ from a point of $\sigma$ next to $v$, with $v$ at the origin
and $\sigma$ along $e_1$. The separatrix stops on $\tau$ at a right angle, so
the return is a segment of length $\delta$ along $\pm e_2$. Homotopic to the
curve of a turn, $\lambda$ has a translation as holonomy, which maps $v$
developed at the start onto $v$ developed at the end, at distance $\delta$
from $q$ across $\sigma$. A free homotopy conjugates the holonomy by a motion
whose rotation carries the walking direction of $\gamma$ onto that of
$\sigma$, here the identity.
\end{proof}

\begin{figure}[t]
\centering
% A separatrix that winds round a cylinder, and the turn it keeps shifted
% (Proposition 8, the mechanism of turn 54 of way 0 on the male body).
% \input inside a figure; needs \usepackage{tikz} and
% \usetikzlibrary{arrows.meta,decorations.pathreplacing}.
%
% (b) is the upper part of the band of dislocation.tex (b), the same T-mesh
% (check.py checks its numbers):
% cut open along the trace tau of the vertex v of valence five, the right side
% glued to the left side two units higher, (10,y) ~ (0,y+2) for y > 3; v is at
% (10,3) and (0,5) (and at (0,3), out of the picture with the lower band),
% the vertex c_3 of valence three at (1,4), its corner glued by a quarter-turn
% (see dislocation.tex). Traces: tau = the side x = 10 from v upwards (and its
% copy x = 0 from v upwards); sigma = the separatrix of v that leaves it to the
% right along y = 5; the separatrix of c_3 along y = 4, which stops on tau at
% n = (10,4); the separatrix of c_3 upwards, which stops on sigma at (1,5).
%   sigma goes once round, 10 units, and stops on tau at q = (10,5), at the
%   distance delta = l_e + l_f = 1 + 1 = 2 from v along tau, with
%   e = [v, n] and f = [n, q].
%   The turn gamma at o = 5.5, just above sigma, crosses tau at (10,5.5) =
%   (0,7.5): it comes back 2 higher, s_gamma = 2 = delta, length 10 = the
%   length of sigma (the along component of the holonomy of the loop lambda:
%   sigma from v to q, then tau back to v). gamma is a ring only if delta = 0,
%   that is if sigma comes back onto v.
%   By the balances of the two patches between y = 3, 4 and 5 (right of
%   x = 1), l_e equals the arc of c_3 downwards and l_f its arc upwards: if a
%   bound of separation holds either at length at least one, gamma is a helix
%   under every quantization meeting the bounds.
%   Parity (Proposition 5): 10 = 2 = 0 mod 2.
% (a) is a sketch of the same band on the cylinder, not to scale: sigma and
% gamma rise by delta = 2 rows in one round (the twist of the band above v).
\begin{tikzpicture}[x=0.46cm, y=0.46cm, font=\footnotesize,
    side/.style={draw=black, line width=0.5pt},
    quadgrid/.style={draw=black!18, line width=0.25pt, step={(1,1)}},
    surface/.style={fill=black!3},
    node dot/.style={circle, fill=black, inner sep=0pt, minimum size=2.6pt},
    separatrix/.style={draw=valfive, line width=1.1pt},
    gap/.style={draw=measured, line width=2pt},
    turn/.style={line width=1pt, -{Stealth[length=5pt]}},
    brace/.style={decorate, decoration={brace, amplitude=3pt}, line width=0.5pt},
    val3/.pic={\fill[valthree, draw=black, line width=0.3pt]
      (90:3.4pt) -- (210:3.4pt) -- (330:3.4pt) -- cycle;},
    val5/.pic={\fill[valfive, draw=black, line width=0.3pt]
      (90:3.4pt) -- (162:3.4pt) -- (234:3.4pt) -- (306:3.4pt) -- (18:3.4pt) -- cycle;}]
  \definecolor{loopspiral}{HTML}{D55E00}
  \definecolor{valthree}{HTML}{E8890C}
  \definecolor{valfive}{HTML}{7C3AAD}
  \definecolor{measured}{HTML}{14934A}

  % ---------------------------------------------------------------- (a)
  % The cylinder; a point (x, y) of the band above v sits at the angle
  % 90 - 36 x degrees (v and tau in front) and the height y + 2x/10 - 5.
  \begin{scope}[shift={(-7.5,1.4)}]
    \pgfmathsetmacro{\R}{2.4}
    \pgfmathsetmacro{\k}{0.22}
    \pgfmathsetmacro{\H}{6.6}
    % Back halves of sigma and gamma (x from 2.5 to 7.5), through the cylinder.
    \draw[valfive!70, dashed, line width=0.7pt]
      plot[domain=2.5:7.5, samples=40, variable=\x]
        ({\R*cos(90-36*\x)}, {1.6 + 0.2*\x - \k*\R*sin(90-36*\x)});
    \draw[loopspiral!70, dashed, line width=0.7pt]
      plot[domain=2.5:7.5, samples=40, variable=\x]
        ({\R*cos(90-36*\x)}, {2.1 + 0.2*\x - \k*\R*sin(90-36*\x)});
    % The cylinder, translucent.
    \fill[black!25, fill opacity=0.18] (-\R,0) -- (-\R,\H) arc[start angle=180, end angle=360,
      x radius=\R, y radius=\k*\R] -- (\R,0) arc[start angle=0, end angle=-180,
      x radius=\R, y radius=\k*\R] -- cycle;
    \draw[line width=0.5pt] (-\R,0) -- (-\R,\H) (\R,0) -- (\R,\H);
    \draw[line width=0.5pt] (0,\H) ellipse[x radius=\R, y radius=\k*\R];
    \draw[line width=0.5pt] (\R,0) arc[start angle=0, end angle=-180,
      x radius=\R, y radius=\k*\R];
    \draw[line width=0.35pt, black!40, dashed] (\R,0) arc[start angle=0, end angle=180,
      x radius=\R, y radius=\k*\R];
    % tau, in front, from v upwards.
    \draw[separatrix] (0,1.6-\k*\R) -- (0,\H-\k*\R);
    % Front halves of sigma (from v) and gamma.
    \draw[separatrix]
      plot[domain=0:2.5, samples=20, variable=\x]
        ({\R*cos(90-36*\x)}, {1.6 + 0.2*\x - \k*\R*sin(90-36*\x)});
    \draw[separatrix]
      plot[domain=7.5:10, samples=20, variable=\x]
        ({\R*cos(90-36*\x)}, {1.6 + 0.2*\x - \k*\R*sin(90-36*\x)});
    \draw[loopspiral, line width=1pt]
      plot[domain=0:2.5, samples=20, variable=\x]
        ({\R*cos(90-36*\x)}, {2.1 + 0.2*\x - \k*\R*sin(90-36*\x)});
    \draw[loopspiral, line width=1pt, -{Stealth[length=5pt]}]
      plot[domain=7.5:9.9, samples=20, variable=\x]
        ({\R*cos(90-36*\x)}, {2.1 + 0.2*\x - \k*\R*sin(90-36*\x)});
    % The other separatrices of v, started, and c_3 next to it.
    \foreach \a in {162, 234, 306}
      \draw[valfive!55, line width=0.8pt] (0,1.6-\k*\R) -- ++(\a:0.55);
    \pic at ({\R*cos(54)}, {0.8 - \k*\R*sin(54)}) {val3};
    % The gap delta on tau, from v to q.
    \draw[gap] (0,1.6-\k*\R) -- (0,3.6-\k*\R);
    \pic at (0,1.6-\k*\R) {val5};
    \node[node dot] at (0,3.6-\k*\R) {};
    \node[left=3pt] at (0,1.6-\k*\R) {$v$};
    \node[right=1pt] at (0,3.6-\k*\R) {$q$};
    \node[valfive, left] at (0,4.8-\k*\R) {$\tau$};
    \node[valfive, right] at (\R,1.95) {$\sigma$};
    \node[loopspiral, right] at (\R,2.75) {$\gamma$};
    \node[anchor=north] at (0,-2.0) {\small (a) the cylinder};
  \end{scope}

  % ---------------------------------------------------------------- (b)
  \begin{scope}
    \fill[surface] (0,1.6) rectangle (10,3) (1,3) rectangle (10,4) (0,4) rectangle (10,8.6);
    \draw[quadgrid] (0,1.6) grid (10,3);
    \draw[quadgrid] (1,3) grid (10,4);
    \draw[quadgrid] (0,4) grid (10,8.6);
    % Traces: the loop of v at y = 3, the separatrices of c_3, and the cut
    % edges of the corner of c_3; then tau and sigma.
    \draw[side] (0,1.6) -- (0,3) (10,1.6) -- (10,3);
    \draw[side] (0,3) -- (10,3) (1,4) -- (10,4) (1,4) -- (1,5);
    \draw[side] (1,4) -- (1,3) (1,4) -- (0,4) -- (0,5);
    \draw[separatrix] (10,3) -- (10,8.6) (0,5) -- (0,8.6);
    \draw[separatrix] (0,5) -- (10,5);
    \draw[gap] (10,3) -- (10,5);
    \foreach \p in {(1,3), (0,4), (10,4), (0,6), (1,5), (10,5), (0,7)}
      \node[node dot] at \p {};
    % The turn gamma, and where it comes back.
    \draw[turn, loopspiral] (0,5.5) -- (10.1,5.5);
    \draw[loopspiral, line width=0.5pt, fill=white] (0,5.5) circle (1.8pt);
    \draw[loopspiral, dashed, line width=0.8pt, -{Stealth[length=4pt]}] (0,7.5) -- (2,7.5);
    \draw[loopspiral, line width=0.5pt, fill=white] (0,7.5) circle (1.8pt);
    \node[above, loopspiral] at (5,5.5) {$\gamma$};
    \draw[brace] (-0.3,5.5) -- (-0.3,7.5) node[midway, left=3pt] {$s_\gamma = 2$};
    % Singular vertices and names.
    \pic at (1,4) {val3};
    \foreach \p in {(0,3), (10,3), (0,5)} \pic at \p {val5};
    \node[above right=-1pt] at (1,4) {$c_3$};
    \node[left=1pt] at (0,4.9) {$v$};
    \node[right=1pt] at (10,3) {$v$};
    \node[right=1pt] at (10,4) {$n$};
    \node[right=1pt] at (10,5) {$q$};
    \node[valfive, below] at (5,5) {$\sigma$};
    \node[valfive, right] at (10,8) {$\tau$};
    \node[measured, left] at (10,3.5) {$e$};
    \node[measured, left] at (10,4.5) {$f$};
    \draw[brace, decoration={mirror}] (11.3,3) -- (11.3,5)
      node[midway, right=3pt] {$\delta = \ell_e + \ell_f = 2$};
  \end{scope}
  \node[anchor=north] at (5,0.9) {$s_\gamma = \delta = 2$: a ring only if $\sigma$ comes back onto $v$};
  \node[anchor=north] at (5,-0.6) {\small (b) cut open along $\tau$};
\end{tikzpicture}

\caption{A separatrix that winds (Proposition~\ref{prop:winding}). (a)~On
a cylinder such as a wrist, the separatrix $\sigma$ of the vertex $v$ of
valence five goes once round and stops on another separatrix $\tau$ of $v$.
(b)~Cut open along $\tau$: $\sigma$, of length $10$, stops at $q$, at
$\delta = \ell_e + \ell_f = 2$ from $v$, and the turn $\gamma$ beside
$\sigma$ comes back $2 = \delta$ rows higher; it is a ring only if $\sigma$
comes back onto $v$. At the left wrist of the male body $\delta = \ell_{1920}
+ \ell_{1921}$, kept positive by a bound of separation.}
\label{fig:separatrix}
\end{figure}

At the left wrist the turn ($42$ of way 0, length $6.99$, real shift
$-0.305$) bounds a cap of four vertices of valence three, whose Burgers sum
gives its shift as a form without sign. The sign comes from outside: the
separatrix $964$ of the vertex $246$, of valence five, runs $6.994$ edges once
round the wrist and stops on another of its separatrices at $\delta =
\ell_{1920} + \ell_{1921} = 0.305$ from $246$; the shift is $-\delta$, and
$\ell_{1920}$ equals the length of the held arc $1937$, both modulo the
balances: whence the certificate and $\shift \le -1$ at every admissible whole
length. We checked the identity, not the homotopy (measured).

\paragraph{Larger conflicts.}
The three turns at the right eye wind round one 3--5 pair, and their
combination is the length of a staircase of three arcs through the held arc
$619$; the four at the eyes wind round a 3--5 pair and two vertices of valence
three, a dislocation and an exchange in homology. The fractional multipliers of Table~\ref{tab:iis} come
from asking the bounded arcs to weigh one in total; at the left hand, where
the turns cross each other $10$ to $46$ times, the primitive ones give
\[
  -(7\shift_{66} + 5\shift_{71} + 2\shift_{75} + 3\shift_{76}) = 2F =
  10\,\ell_{1805} + 2\ell_{1817} + 4\ell_{1853} + 2(\ell_{1870} + \ell_{1879}
  + \ell_{1969} + \ell_{1998} + \ell_{2029})
\]
modulo the balances. The coefficient $5$ is a dislocation: a separatrix of the
vertex $231$ (valence three) stops on one of $228$ (valence five), and arc
$1805$, $0.063$ edges long, joins that junction to $228$, so $\ell_{1805}$ is
the offset of the pair across the line of the field that nearly joins them,
Proposition~\ref{prop:dislocation}(ii) with $\delta = \ell_{1805}$ (hypothesis
on $\kappa$ not checked); their
twisted windings (Section~\ref{sec:fox}) contribute exactly $10\,\ell_{1805}$
(measured).

\paragraph{Local dislocations.}
The curves of an irreducible set need not bound one planar region round a
cluster of singularities: configurations B to F and the sets of four and seven
turns do not. What the conflicts share is local.

\begin{question}[Are conflicts sums of local dislocations?]\label{conj:local}
On a closed surface of genus zero, let $H$ be an irreducible set of turns that
cannot be closed together within the bounds of separation. Does some
combination $\sum_{\gamma \in H} u_\gamma \shift_\gamma$ with whole
multipliers that certifies the conflict (Remark~\ref{rem:definite}) equal,
modulo the balances, a sum with non-negative whole weights of
\begin{enumerate}
\item[(i)] lengths of staircases (Proposition~\ref{prop:staircase}) between a
  vertex of valence three and one of valence five, or between two vertices of
  equal valence exchanged as in Proposition~\ref{prop:swap}, and
\item[(ii)] shifts of caps: sign-definite forms carried by curves that wind
  once round singular vertices of total index four, as the wrist turn and its
  separatrix (Proposition~\ref{prop:winding}),
\end{enumerate}
one term at least containing an arc held open?
\end{question}

It holds with one term, up to the unchecked homotopies, for the wrist,
configurations A to D and the right eye; Section~\ref{sec:combs} settles it on
the body: no as stated, yes and empty when read loosely.

\section{Any combination of turns}\label{sec:fox}

Winding numbers are the homology class of a weighted cycle $\sum_\gamma
u_\gamma \bar\gamma$, which is commutative, while the holonomy is not; the
free differential calculus of Fox~\cite{fox1953free} keeps the order. Let
$\Sigma$ be closed of genus zero, with singular vertices $c_1, \dots, c_n$
(Proposition~\ref{prop:burgers-any} adds borders and handles). Fix $x_0$ inside
a patch, with its frame, and disjoint paths $\pi_i$ from $x_0$ to $c_i$; let
$\lambda_i$ follow $\pi_i$, go round $c_i$ with $c_i$ on its right, and come
back. The fundamental group of $\Sigma^\circ$ is generated by the $\lambda_i$
with a single relation, their product in the order in which the paths leave
$x_0$. Let $p_i$ be the developed position of $c_i$ at the end of $\pi_i$ and
$R_i = R^{4 - k(c_i)}$, so that $\hat h(\lambda_i) = \rho_i$, $\rho_i(x) =
R_i(x - p_i) + p_i$. Write $\hat h(\gamma)(x) = \rho(\gamma)x + \tau(\gamma)$;
then $\tau(\alpha\beta) = \tau(\alpha) + \rho(\alpha)\,\tau(\beta)$. Let $F$
be the free group on $x_1, \dots, x_n$, mapped onto the fundamental group by
$x_i \mapsto \lambda_i$, with the Fox derivatives $\partial x_j/\partial x_i =
\delta_{ij}$, $\partial(uv)/\partial x_i = \partial u/\partial x_i +
u\,\partial v/\partial x_i$. Extended linearly to the group ring, $\rho$ takes
the values $a_0 I + a_1 R + a_2 R^2 + a_3 R^3$, and the augmentation $\epsilon$
maps such an element to $a_0 + a_1 + a_2 + a_3$. The first identity below is
the defining property of Fox derivatives for a crossed
homomorphism~\cite{fox1953free}, and the $p_i$ are, up to a common motion,
twisted-cohomology coordinates of the flat
metric~\cite{thurston1998shapes,veech1993flat}; what is new is the reading on
combinations of turns.

\begin{proposition}[Fox--Burgers formula]\label{prop:fox}
For every word $w$ of $F$,
\[
  \tau(w) = \sum_{i=1}^n \rho\Bigl(\frac{\partial w}{\partial x_i}\Bigr)
  (I - R_i)\,p_i .
\]
Let $\gamma_1, \dots, \gamma_m$ be turns, $\kappa_r$ a path from $x_0$ to the
start of $\bar\gamma_r$, $f_r$ the direction $e_2$ of $\gamma_r$ carried back
to $x_0$ along $\kappa_r$ (one of $\pm e_1$, $\pm e_2$), and $w_r$ a word for
$\kappa_r \bar\gamma_r \kappa_r^{-1}$. For multipliers $u_r$ define the
\emph{twisted winding} of the weighted cycle round $c_i$,
\[
  W_i = -\sum_{r=1}^m u_r\, \rho\Bigl(\frac{\partial w_r}{\partial x_i}\Bigr)^{\!\top} f_r
  \in \Z^2 \text{ (for whole } u_r) .
\]
Then, for all balanced lengths,
\[
  \sum_{r=1}^m u_r\, \shift_{\gamma_r}(\ell) = \sum_{i=1}^n \bigl\langle W_i,\
  (I - R_i)\,p_i \bigr\rangle ,
  \qquad
  \sum_{i=1}^n (I - R_i)^\top W_i = 0 ,
\]
so the sum does not change when all the $p_i$ are moved by one translation: it
is a combination of separations $p_i - p_j$. The ordinary winding number of
$\sum_r u_r \bar\gamma_r$ round $c_i$ is, up to a constant common to all $i$,
$w_i = \sum_r u_r\, \epsilon(\partial w_r/\partial x_i)$. The words $w_r$ are
defined up to the relation, and the $W_i$ accordingly: inserting a conjugate
$g\varrho g^{-1}$ of the relation $\varrho$ in $w_r$, after a prefix $a$ of
$w_r$, subtracts $u_r \rho(a g\, \partial \varrho/\partial x_i)^\top f_r$ from
every $W_i$ (with $a$ empty, $u_r \rho(g\, \partial \varrho/\partial x_i)^\top f_r$),
and leaves the sum unchanged.
\end{proposition}

\begin{proof}
Both $w \mapsto \tau(w)$ and the right-hand side satisfy $\phi(uv) = \phi(u) +
\rho(u)\phi(v)$, and they agree on the letters and the empty word, so they
agree everywhere by induction on the length (for $x_i^{-1}$, use $0 =
\phi(x_i x_i^{-1})$). For a turn, $\hat h(\kappa_r \bar\gamma_r
\kappa_r^{-1})$ is the translation by $M_r T_r$, with $T_r$ the translation of
$\bar\gamma_r$ in its own frame and $M_r$ the rotation of $\hat h(\kappa_r)$;
so $\shift_{\gamma_r} = -\langle M_r e_2, M_r T_r \rangle = -\langle f_r,
\tau(w_r) \rangle$, and $\langle f, \rho(D) b\rangle = \langle \rho(D)^\top f,
b \rangle$ gives the formula. The fundamental formula $\sum_i (\partial
w/\partial x_i)(x_i - 1) = w - 1$ maps under $\rho$ to $\sum_i \rho(\partial
w_r/\partial x_i)(R_i - I) = \rho(w_r) - I = 0$, the rotation of a turn being
trivial; transposed and summed with the $u_r f_r$, it is the second identity.
Finally $\epsilon(\partial w/\partial x_i)$ is the exponent sum of $x_i$ in
$w$, the homology class of $w$ on the basis of the $\lambda_i$, defined up to
the relation, in which each $\lambda_i$ appears once.
\end{proof}

For a simple turn $\rho(\partial w/\partial x_i) = Q_i$ and the formula is
Proposition~\ref{prop:burgers}(a). In general a cycle may go round $c_i$
several times after different rotations; modulo $2$ these disappear.

\begin{corollary}[Checkerboard rule for any combination]\label{cor:checker}
On a closed surface of genus zero, let $\ell$ be whole and balanced and the
multipliers $u_r$ whole. Colour each singular vertex $c_i$ by the parity
$\chi_i$ of $p_{i,1} + p_{i,2}$; it does not depend on the paths $\pi_i$, and
for positive lengths it is the colour of $c_i$ in the bipartition of the quad
mesh. Then
\[
  \sum_r u_r\, \shift_{\gamma_r}(\ell) \equiv \sum_r u_r\, \alpha_{\gamma_r}(\ell)
  \equiv \sum_{i:\ k(c_i) \text{ odd}} w_i\, \chi_i \pmod 2 ,
\]
with $w_i$ the winding numbers of $\sum_r u_r \bar\gamma_r$, whatever the
common constant and whatever the orientation of each $\bar\gamma_r$. In
particular, a combination whose cycle winds an even number of times round every
vertex of odd valence takes even values at every whole length.
\end{corollary}

\begin{proof}
Since $(R - I)^2 = -2R$, $R^k(I - R^{\pm1}) \equiv I - R^{\pm1}$ modulo $2$, so
$\rho(D)(I - R_i) \equiv \epsilon(D)(I - R_i)$ for $k(c_i)$ odd, and $\equiv 0$
for $k(c_i)$ even, where $I - R_i \in \{0, 2I\}$; and $\langle f, (I -
R^{\pm1})p\rangle \equiv p_1 + p_2$ for $f \in \{\pm e_1, \pm e_2\}$. This
gives the sum of Proposition~\ref{prop:fox} modulo $2$, and the direction
$e_1$ of each turn that of $\sum_r u_r \alpha_{\gamma_r}$. Rotations about
whole points and quarter-turns keep $p_1 + p_2 \bmod 2$, so $\chi_i$ does not
depend on the paths. A whole translation of the origin flips all colours and
adds $\sum_{k(c_i)\text{ odd}} w_i$, even because the rotation of the cycle is
trivial ($\sum_i w_i (4 - k(c_i)) \equiv 0 \pmod 4$); adding $q$ to all $w_i$
adds $q \sum_{k(c_i)\text{ odd}} \chi_i$, the value modulo $2$ of the
relation, a loop with zero translation, hence even. Reversing a turn changes
signs only. For positive lengths the developed parity colours the quad graph
properly, as in Proposition~\ref{prop:parity}.
\end{proof}

The corollary contains Propositions~\ref{prop:checkerboard},
\ref{prop:dislocation}(iii) and~\ref{prop:parity}. As for crystal
defects~\cite{kleman2008disclinations} and in~\cite[\S7]{izmestiev2013torus},
winding numbers do not determine the holonomy:

\begin{proposition}[Winding numbers do not suffice]\label{prop:commutator}
Let $c_a$, $c_b$ be singular vertices of odd valence developed at distinct
points $p_a \neq p_b$. The loop $\lambda_a \lambda_b \lambda_a^{-1}
\lambda_b^{-1}$ winds zero times round every singular vertex, and its holonomy
is the translation by $(I - R_a)(I - R_b)(p_a - p_b)$, of length $2|p_a - p_b|
\neq 0$. No function of the winding numbers gives the holonomy of the loops
with trivial rotation.
\end{proposition}

\begin{proof}
For $w = x_a x_b x_a^{-1} x_b^{-1}$, $\rho(\partial w/\partial x_a) = \rho(1 -
x_a x_b x_a^{-1}) = I - R_b$ and $\rho(\partial w/\partial x_b) = \rho(x_a -
w) = R_a - I$, the rotations commuting; Proposition~\ref{prop:fox} gives the
translation, and for odd valences $I - R_a$, $I - R_b$ are $\sqrt2$ times
rotations. The exponent sums of $w$ are zero.
\end{proof}

\paragraph{On the male body.}
Cut into one chart along a spanning tree of the nodes, the T-mesh gives every
turn a word, and Proposition~\ref{prop:fox} its translation, for all
$1{,}299$ turns. Corollary~\ref{cor:checker} holds on all of them and on the
first $14$ sets of Table~\ref{tab:iis}, fractional multipliers scaled to whole
ones (the set of $15$ turns was not tried); at the left hand the
cycle winds an even number of times round every vertex of odd valence, which
is why $-2F$ is even. Four turns satisfy $w_{18} = w_{22} + w_{101} + w_{65}$
round every singular vertex while their translations do not, in any
quarter-turn frames: a ``winding Burgers formula'' with numbers for weights is
false (measured).

\subsection{Combs}\label{sec:combs}

The twisted windings depend on the chart; combs do not. In the T-meshes of
openclare, traced from the separatrices of the singular vertices only, every
trace of a closed surface starts at a singular vertex (feature lines or seeded
traces would break this; we assume it). For an arc $a$, along $e_a$ from the node $n_0$
to the node $n_1$, the \emph{foot} of $n_j$ is the node itself if it is
singular, and otherwise the singular vertex where the trace through $n_j$
perpendicular to $a$ starts. We assume this vertex unique: a single trace
perpendicular to $a$ meets $n_j$ (not two, from opposite sides), and it is neither a separatrix
connection nor a loop; otherwise a foot must be chosen, and different choices
give different dislocations. The \emph{comb} $\sigma_a$ of $a$ runs from the
foot $c_0$ of $n_0$ along its trace to $n_0$, along $a$ to $n_1$, and along the
other trace to the foot $c_1$ of $n_1$, through T-junctions only
(Figure~\ref{fig:comb}). The \emph{graph of feet} has the singular vertices of
odd valence as vertices and the arcs whose two feet are distinct vertices of
odd valence as edges.

\begin{figure}[t]
\centering
% The comb of an arc (Section "Combs" and Lemma "Combs").
% \input inside a figure; needs \usepackage{tikz} and
% \usetikzlibrary{arrows.meta,decorations.pathreplacing}.
%
% The T-mesh is that of dislocation.tex (b) and separatrix.tex (b): a band cut
% open along the trace tau of the vertex v of valence five, the right side
% glued to the left side two units higher above y = 3 and at the same height
% below; v is drawn at (0,3), (10,3) and (0,5), the vertex c_3 of valence three
% at (1,4), and the empty corner next to c_3 is glued by a quarter-turn about
% c_3 (see dislocation.tex). Every node is at a whole position; check.py
% checks the numbers below.
%
% The arc a is the piece of the separatrix sigma of v from the T-junction
% n_0 = m = (1,5), where the separatrix of c_3 upwards stops on sigma, to the
% T-junction n_1 = q = (10,5), where sigma stops on tau; e_a = (1,0), l_a = 9.
% Foot of n_0: the trace through m perpendicular to a is the separatrix of c_3
% upwards, which starts at c_3: c_0 = c_3. Foot of n_1: the trace through q
% perpendicular to a is tau, which starts at v: c_1 = v. The comb sigma_a runs
% from c_3 up to m (1 unit), along a to q (9 units), and down tau through the
% T-junction n = (10,4) to v (2 units): it passes through T-junctions only.
% Developed along it, p_{c_0} = (1,4) and p_{c_1} = (10,3), and
%   < e_a, p_{c_1} - p_{c_0} > = < (1,0), (9,-1) > = 9 = l_a
% (Lemma "Combs"); the teeth only add the across component -1.
% The feet are distinct, of valences 3 and 5: the comb reads the length of a
% as that of a dislocation (Proposition 3); with two feet of valence 3 it would
% be an exchange (Proposition 7), with a single foot a winding separatrix
% (Proposition 8).
\begin{tikzpicture}[x=0.5cm, y=0.5cm, font=\footnotesize,
    side/.style={draw=black, line width=0.5pt},
    quadgrid/.style={draw=black!18, line width=0.25pt, step={(1,1)}},
    surface/.style={fill=black!3},
    node dot/.style={circle, fill=black, inner sep=0pt, minimum size=2.6pt},
    comb/.style={draw=measured, line width=1.8pt, line join=round},
    arc mark/.style={line width=2.2pt, draw=black},
    dimension/.style={line width=0.5pt, {Bar[width=5pt]}-{Bar[width=5pt]}},
    val3/.pic={\fill[valthree, draw=black, line width=0.3pt]
      (90:3.4pt) -- (210:3.4pt) -- (330:3.4pt) -- cycle;},
    val5/.pic={\fill[valfive, draw=black, line width=0.3pt]
      (90:3.4pt) -- (162:3.4pt) -- (234:3.4pt) -- (306:3.4pt) -- (18:3.4pt) -- cycle;}]
  \definecolor{valthree}{HTML}{E8890C}
  \definecolor{valfive}{HTML}{7C3AAD}
  \definecolor{measured}{HTML}{14934A}

  \begin{scope}
    \fill[surface] (0,1.6) rectangle (10,3) (1,3) rectangle (10,4) (0,4) rectangle (10,8);
    \draw[quadgrid] (0,1.6) grid (10,3);
    \draw[quadgrid] (1,3) grid (10,4);
    \draw[quadgrid] (0,4) grid (10,8);
    % All traces of the T-mesh in the picture.
    \draw[side] (0,1.6) -- (0,3) (10,1.6) -- (10,8) (0,5) -- (0,8);
    \draw[side] (0,3) -- (10,3) (1,4) -- (10,4) (1,4) -- (1,5) (0,5) -- (10,5);
    \draw[side] (1,4) -- (1,3) (1,4) -- (0,4) -- (0,5);
  \end{scope}
  % The comb of a: from its foot c_3 up to m, along a, down tau to its foot v.
  \draw[comb] (1,4) -- (1,5) -- (10,5) -- (10,3);
  \draw[comb, line width=3pt] (1,5) -- (10,5);
  \foreach \p in {(1,3), (0,4), (10,4), (0,6), (1,5), (10,5), (0,7)}
    \node[node dot] at \p {};
  \pic at (1,4) {val3};
  \foreach \p in {(0,3), (10,3), (0,5)} \pic at \p {val5};
  % Names.
  \node[above] at (5.5,5) {$a$};
  \draw[-{Stealth[length=4pt]}, line width=0.5pt] (6.6,5.75) -- (8,5.75)
    node[right, inner sep=1pt] {$e_a$};
  \node[above left=-1pt] at (1,5) {$n_0$};
  \node[above right=-1pt] at (10,5) {$n_1$};
  \node[below right=-1pt] at (1,4) {$c_0 = c_3$};
  \node[right=2pt] at (10,3) {$c_1 = v$};
  \node[left=1pt] at (0,4.95) {$v$};
  \node[left=1pt] at (0,3) {$v$};
  \node[measured, right=1pt] at (10,4.5) {$\sigma_a$};
  \node[right=1pt] at (10,7.4) {$\tau$};
  % The developed feet and the e_a component of their separation.
  \draw[dimension] (1,2.2) -- (10,2.2) node[midway, fill=black!3, inner sep=1.5pt]
    {$\langle e_a,\, p_{c_1} - p_{c_0}\rangle = \ell_a = 9$};
  \draw[black!50, dotted, line width=0.6pt] (1,2.2) -- (1,4) (10,2.2) -- (10,3);
\end{tikzpicture}

\caption{The comb $\sigma_a$ of an arc $a$ (Lemma~\ref{lem:comb}), $\ell_a =
9$: from the foot $c_0$ of one end of $a$, of valence three, up its trace,
along $a$, and down the other perpendicular trace to the foot $c_1$, of
valence five. Every step but $a$ is orthogonal to $a$, so $\langle e_a,
p_{c_1} - p_{c_0} \rangle = 9 = \ell_a$. This T-mesh is that of
Figures~\ref{fig:dislocation}(b) and~\ref{fig:separatrix}(b): the staircase,
the winding separatrix and the comb are three readings of the same shift
$2$.}
\label{fig:comb}
\end{figure}

\begin{lemma}[Combs]\label{lem:comb}
Let $p_{c_0}$, $p_{c_1}$ be the feet of an arc $a$ developed along its comb.
Then, for all balanced lengths,
\[
  \langle e_a,\ p_{c_1} - p_{c_0} \rangle = \ell_a .
\]
If $c_0 \neq c_1$ have odd valence, $2\ell_a$ is, up to sign, the form of the
twisted class $\lambda_{c_0}\lambda_{c_1}^{\pm1} \otimes v$ of the loops round
the two feet reached along the comb, for a whole $v$ with $(I - R^{\pm1})^\top v
= 2e_a$: the pair is a dislocation if the valences differ
(Proposition~\ref{prop:dislocation}) and an exchange if they are equal
(Proposition~\ref{prop:swap}). If $c_0 = c_1$, $\ell_a$ is the shift of a
winding separatrix (Proposition~\ref{prop:winding}).
\end{lemma}

\begin{proof}
Every step of the comb but $a$ runs along a trace perpendicular to $a$, through
no cone point, and adds a vector orthogonal to $e_a$; the step along $a$ adds
$\ell_a e_a$. The loop $\lambda_{c_0} \lambda_{c_1}^{\pm1}$, with the sign that
makes its rotation trivial, has the translation $\pm(I - R^{\pm1})(p_{c_0} -
p_{c_1})$ as in Propositions~\ref{prop:dislocation} and~\ref{prop:swap};
pairing it with $v$ gives $\pm 2\ell_a$, and $v$ is whole because $I -
R^{\pm1}$ maps $\Z^2$ onto the vectors whose coordinates have equal parity,
among them $2e_a$. If $c_0 = c_1$, the second trace is a separatrix of the same
vertex stopping on the first at the distance $\ell_a$.
\end{proof}

A \emph{closed walk of combs} is a sequence of arcs $a_1, \dots, a_k$ whose
feet are $c_{i-1}$ and $c_i$, of odd valence, with $c_k = c_0$; its
\emph{reading} is $\sum_i \ell_{a_i}$. A walk of length one is a winding
separatrix, one of length two between valences three and five the staircase of
Proposition~\ref{prop:staircase}, and one round four vertices of valence three
a cap.

\begin{proposition}[Certificates are sums of combs]\label{thm:combs}
Let $c = \sum_a c_a\ell_a$, $c_a \in \Z_{\ge0}$, be a certificate
(Remark~\ref{rem:definite}) whose arcs have feet of odd valence and whose odd
part, the arcs with $c_a$ odd, is a cycle modulo~$2$ of the graph of feet.
Then $c$ is a sum with whole non-negative weights of comb readings $2\ell_a$,
each the form of a dislocation or an exchange (Lemma~\ref{lem:comb}), and of
readings of closed walks of combs. With real weights, every non-negative form
supported on arcs whose two feet are of odd valence, distinct or equal, is a
non-negative combination of comb readings $2\ell_a$ and of winding separatrices
$\ell_a$, without condition of parity.
\end{proposition}

\begin{proof}
The odd part splits into closed walks along its edges (Veblen), and what
remains is $\sum_a \lfloor c_a/2 \rfloor\,2\ell_a$. With real weights, $c =
\sum_a \tfrac{c_a}{2}\,(2\ell_a)$ on the arcs with two distinct feet, and
$c_a\ell_a$ on those whose feet are one vertex.
\end{proof}

The proof is elementary: the content is in the hypothesis. That the odd part of
a certificate is a cycle of the graph of feet is observed on all $25$
certificates of the body (exact) but not proved; Corollary~\ref{cor:checker}
does not give it, since it controls the parity of $p_{i,1} + p_{i,2}$ while a
comb reads one coordinate (in Figure~\ref{fig:comb} both feet have the same
colour and $\ell_a = 9$).

\begin{remark}[Closed walks as twisted cycles]
We expect the reading of a closed walk of combs, $\sum_i \langle e_{a_i} -
e_{a_{i+1}},\ p_{c_i}\rangle$ by a summation by parts, to be the Burgers form
of a twisted cycle near its combs; but the two combs through $c_i$ develop it
in charts that differ by a rotation, and a proof must develop the whole walk
in one chart. We use closed walks only as forms.
\end{remark}

\begin{proposition}[The Burgers pairing]\label{prop:pairing}
Let the surface be closed of genus zero, with no singular vertex of valence
divisible by four, and let every trace of the T-mesh start at a singular
vertex. Then the map that sends a class of the twisted homology
$H_1(\Sigma^\circ; \R^2_\rho)$ to its form $\sum_i \langle W_i, (I - R_i)p_i
\rangle$ on the balanced lengths is onto the dual of the balanced lengths. Its
kernel contains the class across each separatrix connection. If some singular
vertices have valence divisible by four, no form depends on their developed
positions, and the map need not be onto.
\end{proposition}

\begin{proof}
If every such form vanishes at balanced $\ell$, every loop of $\Sigma^\circ$
with trivial rotation has a holonomy without translation
(Proposition~\ref{prop:fox}). The commutator of $\lambda_a$ and $\lambda_b$
has trivial rotation and translation $(I - R_a)(I - R_b)(p_a - p_b)$
(Proposition~\ref{prop:commutator}); no $R_i$ is the identity, so $I - R_i$ is
invertible ($\sqrt2$ times a rotation, or $2I$), and every singular vertex
develops at one point $q$. As the $\rho_i$ fix $q$ and generate the holonomy on a closed
sphere, every singular vertex develops at $q$ along any path; every arc has
feet, so every comb offset of Lemma~\ref{lem:comb} is zero, and $\ell = 0$.
The image has no common zero but $0$. Across a connection the
separation of its two vertices is zero across the trace, and the twisted class
pairing the across direction with it has the zero form. A vertex of valence
divisible by four has $(I - R_i)p_i = 0$: its position is a relative period,
which the lengths can move without changing any holonomy.
\end{proof}

\paragraph{On the male body.}
Lemma~\ref{lem:comb} holds on all $2{,}069$ arcs, $2{,}046$ of which have two
distinct feet of odd valence: nearly every arc is a local dislocation or
exchange. On this T-mesh (way~1) each node meets a single perpendicular trace,
and one arc, $1676$, ends on the separatrix connection, which joins the
vertices $204$ and $210$; we took as its foot the vertex the connection was
traced from, $204$, and no certificate of way~1 contains the arc. The
twisted homology has dimension $518 = 2(n - 2)$ and the pairing rank $517$,
its kernel the class across the one separatrix connection; we conjecture that
without vertices of valence divisible by four the kernel is spanned by the
classes across the connections. All $25$ certificates of the conflicts
examined have a representative whose odd part is a cycle of the graph of feet
(exact). The wrist is a comb with both feet at $246$, the winding separatrix;
A, C and D are closed walks of two combs (A the staircase $\ell_{611} +
\ell_{624}$), B of four, the right eye of three, E of six, of total index
$-2$, and F of four round a cap; the left hand is five times the dislocation
$\ell_{1805}$ plus eight comb readings on a tree of seven arcs, and the set of
seven of way~7 a tree of seven combs (measured).

\paragraph{Question~\ref{conj:local}, read loosely, is true and empty.}
Allowing every comb reading and every closed walk of combs,
Proposition~\ref{thm:combs} answers it for every certificate whose odd part is
a cycle of the graph of feet, all those of the body; but this follows from
Farkas' lemma and a parity, and says nothing about which configurations
conflict. \emph{Read strictly, it is false}: restricted to staircases of
weight one, monotone walks round a cap and winding separatrices, exact
rational balanced lengths make every such form non-negative and the
certificates of E and F negative, $-2.9967$ and $-1.9973$ (checked in integer
arithmetic). These forms are infinitely many, but each is the length of a walk
on a graph whose steps carry the arc lengths, so a shortest-path search finds
a negative one if any. In floating point the same test also separates the four
turns at the eyes and puts the left hand inside the cone (measured). Whether
every certificate splits into combs on disjoint sets of vertices, each of its
sign, is open: the trees of the left hand and of way~7 share vertices.

\section{Borders and handles}\label{sec:beyond}

Models with borders (the eye sockets of Suzanne, the mouth of Jerry) and
handles (Bob, a torus; part and cage) need the forms below, measured on the
T-meshes of these bench models with an engine version not recorded. A border
node $b$ is a corner of angle $a(b)\pi/2$, of
\emph{index} $2 - a(b)$; for a border component $\beta_j$, $L_j(\ell)$ is the
sum of the lengths of its arcs.

\begin{lemma}[Gauss--Bonnet with borders]\label{lem:gauss-border}
Let $D$ be a region bounded by curves of turns and by border components. Then
\[
  \sum_{c \in D} (4 - k(c)) + \sum_{b \in D \cap \partial\Sigma} (2 - a(b))
  = 4\chi(D) - \Theta,
\]
where $\Theta$ is the turning of the curves of the turns at their corners, in
quarter-turns: zero for the side of a simple curve.
\end{lemma}
\begin{proof}
Gauss--Bonnet for a flat surface with cone points and a piecewise straight
boundary, border nodes being corners of turning $(2 - a)\pi/2$.
\end{proof}

A side of a simple turn holding $b$ border components without corner and $h$
handles has index $4(1 - b - 2h)$, and may hold no singular vertex at all; the
lemma holds on the $1{,}076$ sides of turns of Suzanne and Jerry (measured).
For positive lengths the homomorphism of the next theorem is the cycle parity
of the quad graph, a $\Z/2$-homology invariant of any
quadrangulation~\cite[Lemma~3.1]{kaiser2015colouring}
extended here to turns at all balanced
lengths.

\begin{theorem}[Parity of a turn on any surface]\label{thm:parity-any}
Let $\ell$ be whole and balanced. There is a homomorphism $\Omega_\ell\colon
H_1(\Sigma; \Z/2) \to \Z/2$ such that for every turn
\[
  \alpha_\gamma(\ell) + \shift_\gamma(\ell) \equiv \Omega_\ell([\bar\gamma]) \pmod 2 .
\]
Its value on a border component is the parity of $L_j(\ell)$, the sum of all
$L_j(\ell)$ is even, and for positive lengths $\Omega_\ell(z)$ is the parity of
the length of any cycle of the quad graph in the class $z$; $\Omega_\ell = 0$
exactly when the quad graph is bipartite. On a surface of genus zero with
borders, $\alpha_\gamma + \shift_\gamma$ has the parity of the total length of
the border components on either side of $\bar\gamma$, or of the number of times
$\bar\gamma$ crosses any system of disjoint arcs joining the border components
of odd length in pairs.
\end{theorem}
\begin{proof}
As in Proposition~\ref{prop:parity}, $\gamma \mapsto \pi(\hat h(\gamma))$ is a
homomorphism, zero on the loops round the singular vertices (quarter-turns
about whole points); with values of exponent two it factors through
$H_1(\Sigma; \Z/2)$, those loops spanning the kernel of $H_1(\Sigma^\circ)
\to H_1(\Sigma)$. For a turn $\hat h(\bar\gamma)$ is $(\alpha_\gamma,
\pm\shift_\gamma)$ up to a quarter-turn; a border component comes back
displaced by $\sum_a \pm\ell_a v_a$, of the parity of $L_j(\ell)$, and the
components add up to the boundary, zero in homology. Along a cycle of the quad
graph every edge changes the colour. In genus zero the border components span
$H_1(\Sigma; \Z/2)$ and $\bar\gamma$ is homologous to those on one side; the
last form is Lefschetz duality.
\end{proof}

On Suzanne, whose eye sockets are borders of $23$ and $25$ edges, the turn
round an eye socket is a ring of $23$ quads at the lengths of one way and of
$24$ at another, against Proposition~\ref{prop:parity}; its shift is zero at
every balanced length:

\begin{proposition}[Collars]\label{prop:collar}
Let the curve of a turn $\gamma$ and a border component $\beta_j$ without
corner bound an annulus holding no singular vertex. Then $\shift_\gamma(\ell)
= 0$ and $\alpha_\gamma(\ell) = L_j(\ell)$ for all balanced lengths.
\end{proposition}
\begin{proof}
The holonomy of $\beta_j$ is the translation by $L_j$ along an axis of the
field. Cut along a path from $\bar\gamma$ to $\beta_j$, the annulus is a disc
without cone point, so the translation of $\bar\gamma$ is $L_j u$ with $u \in
\{\pm e_1, \pm e_2\}$, as in Proposition~\ref{prop:burgers}(b). For positive
lengths $\alpha_\gamma > 0$, so $u = \pm e_1$; the identities are linear and
extend to all balanced lengths.
\end{proof}

The theorem predicts the parity of all $660$ turns of Suzanne, Jerry and a
dome, $61$ of them odd, and of every closed loop of the output meshes ($294$
on the male body); the proposition holds on all $40$ collars; on Bob,
$\Omega_\ell$ depends on the lengths, the three ways giving it $(0, 1)$, $(0,
0)$ and $(1, 0)$ on two handle loops (measured). For combinations,
$\sum_r u_r(\alpha_{\gamma_r} + \shift_{\gamma_r}) \equiv \sum_r u_r\,
\Omega_\ell([\bar\gamma_r])$; the colours of Corollary~\ref{cor:checker} may
depend on the paths, but the count on white points still gives the parity when
every curve lies in a planar region without border.

\begin{proposition}[Burgers sum with borders and handles]\label{prop:burgers-any}
Let the loops $\lambda_i$ round the singular vertices, $\beta_j$ round the
border components and $a_k, b_k$ round the handles generate the fundamental
group of $\Sigma^\circ$. For every word $w$ in them,
\[
  \tau(w) = \sum_i \rho\Bigl(\frac{\partial w}{\partial x_i}\Bigr)(I - R_i)\,p_i
  + \sum_{g \in \{\beta_j, a_k, b_k\}} \rho\Bigl(\frac{\partial w}{\partial g}\Bigr)\tau(g).
\]
For a turn, the shift is the Burgers sum of the singular vertices plus a
\emph{period term}, linear in the lengths and not determined by the developed
positions of the singular vertices. A border component without corner has
holonomy the translation by $L_j$ along its developed direction; a handle
inside the region between two turns contributes the commutator $(I -
\rho(b))\tau(a) - (I - \rho(a))\tau(b)$, zero when both its loops have trivial
rotation.
\end{proposition}
\begin{proof}
Both sides satisfy $\phi(uv) = \phi(u) + \rho(u)\phi(v)$ and agree on the
generators. The commutator of $x \mapsto Ax + a$ and $x \mapsto Bx + b$ with
commuting rotations is the translation by $(I - B)a - (I - A)b$.
\end{proof}

With its border term the sum reproduces all $314$ sides holding a border on
which it was tested, and without it none of the $586$ such sides of Suzanne
and Jerry (measured).

\begin{proposition}[Two crossing turns on a torus]\label{prop:torus}
On a closed surface of genus one, let the curves of two turns be simple and
cross once. Cut along both, the surface is a disc $D$ holding every singular
vertex, and its Burgers sum vanishes: $\sum_i Q_i (I - R_i) p_i = 0$, in the
order of first visit, for all balanced lengths. More generally, if $\kappa$ is
a simple loop crossing the curve of a turn once, with rotation $R^j$, then $(I
- R^j)\,T_\gamma = \pm B_D$: for $j \neq 0$ the translation of the turn is set
by the singular vertices, for $j = 0$ it is a free period and the singular
vertices satisfy $B_D = 0$.
\end{proposition}
\begin{proof}
The border of $D$ is the commutator of the two curves, whose holonomies are
translations and commute; Proposition~\ref{prop:burgers}(a) applies to $D$.
With $\kappa$ of rotation $R^j$, the commutator is the translation by $(I -
R^j)T_\gamma$.
\end{proof}

On Bob, of genus one with ten vertices of valence three and ten of valence
five, the Burgers sum vanishes for all $304$ pairs of simple turns crossing
once, and for none of $90$ shuffled orders (measured). With a single 3--5
pair the two cone points would develop onto one point; that no such torus
exists is the theorem of Barnette, Jucovi\v{c} and
Trenkler~\cite{barnette1971toroidal}, also a consequence of Abel's theorem
($c_5 - c_3$ is principal on an elliptic curve only if $c_5 = c_3$), an
obstruction to existence, not to closing. With $j = 0$ the proposition is the
commutator argument of~\cite[\S7]{izmestiev2013torus}. The conflict met on
Bob is of another kind: $3\shift_{33} - 2\shift_{31}$, for two turns winding
three and two times round the same handle, cancels the period and equals two
arc lengths at a vertex of valence three, one held open (measured).

Proposition~\ref{prop:linear}, Lemma~\ref{lem:gauss} for discs and annuli,
Propositions~\ref{prop:burgers} to~\ref{prop:checkerboard} (for annuli),
\ref{prop:staircase}, \ref{prop:swap}, \ref{prop:winding}
and~\ref{prop:commutator}, Remarks~\ref{rem:definite} and~\ref{rem:exist},
Theorem~\ref{thm:forced} and Corollary~\ref{cor:chain} need no closed sphere;
Lemma~\ref{lem:eight}, Proposition~\ref{prop:parity},
Corollary~\ref{cor:parity}, the count of independent rings, the setting of
Proposition~\ref{prop:fox}, Corollary~\ref{cor:checker} and
Question~\ref{conj:local} do; Appendix~\ref{app:orders} treats triangle
meshes.

\section{Turns every quantization within bounds leaves shifted}\label{sec:forced}

Turned around, the obstructions bound from below the number of turns any
quantization in a set $P$ of lengths leaves shifted.

\begin{theorem}[Forced shifts]\label{thm:forced}
Let $H$ be a finite set of turns of a T-mesh and $P$ a non-empty set of lengths
containing every quantization considered: for instance the balanced lengths
meeting given bounds of separation, $\ell \ge b$, possibly with a cap $\frac1n
\sum_a |\ell_a - t_a| \le D$ on the mean deviation from the targets $t$, or with
a given number of quads. Let
\[
  \tau_P(H) = \min\bigl\{\, |X| : X \subseteq H,\ \text{the turns of } H
  \setminus X \text{ can be closed together by lengths in } P \,\bigr\}.
\]
Then for every $\ell \in P$, whole or not, at least $\tau_P(H)$ turns of $H$
have $\shift_\gamma(\ell) \neq 0$. If $H$ contains $k$ pairwise disjoint subsets
none of which can be closed within $P$, then $\tau_P(H) \ge k$.
\end{theorem}

\begin{proof}
At $\ell \in P$ the turns with zero shift are closed by $\ell$ itself, so the
shifted ones form a set $X$ as in the definition. A set whose removal leaves a
closable set meets each of the $k$ disjoint subsets.
\end{proof}

The theorem is elementary; its content is in what makes subsets infeasible.
It uses neither convexity, nor linearity, nor whole lengths. The minimal
non-closable sets are irreducible infeasible subsystems (IIS), $\tau_P(H)$ is
a minimum IIS cover~\cite{parker1996iis,amaldi2003maxfs}, and the theorem is
the packing bound of that covering problem.

\begin{corollary}[Chains of dislocations]\label{cor:chain}
Let $\gamma_0, \gamma_1, \dots, \gamma_m$ be turns and suppose that for $i = 1,
\dots, m$
\[
  \shift_{\gamma_{i-1}}(\ell) - \shift_{\gamma_i}(\ell) = \varepsilon_i\,
  F_i(\ell) \quad \text{for all balanced } \ell,
\]
with $\varepsilon_i \in \{+1, -1\}$ and forms $F_i$ that are positive on the
whole points of $P$; for instance $F_i$ the length of a staircase
(Proposition~\ref{prop:staircase}) between a vertex of valence three and one of
valence five lying between $\gamma_{i-1}$ and $\gamma_i$
(Propositions~\ref{prop:dislocation} and~\ref{prop:burgers}(b)), the
staircase containing an arc held at length at least one. Let $q$ be the number
of indices $i < m$ with $\varepsilon_{i+1} \neq \varepsilon_i$. Then for every
whole $\ell \in P$ at most $q + 1$ of the turns $\gamma_i$ are rings: at least
$m - q$ are helices. In particular $m + 1$ nested parallel turns separated by
$m$ dislocations whose staircases all climb the same way leave at least $m$
helices.
\end{corollary}

\begin{proof}
If $\gamma_i$ and $\gamma_j$, $i < j$, are rings, $\sum_{k=i+1}^{j}
\varepsilon_k F_k(\ell) = 0$ with $F_k(\ell) > 0$, so a change of sign lies
between $i$ and $j$; between two consecutive rings lies one of its own.
\end{proof}

\begin{remark}[What the bound counts]
The theorem counts turns with a non-zero shift; at lengths where a turn is not
a strip of the mesh~\cite[Prop.~2]{gardent2026loops}, or where one spiral
strip carries several turns of $H$, it is not a count of helices.
\end{remark}

The bound is attained, by the standard correspondence between minimum IIS
covers and maximum feasible subsystems~\cite{amaldi2003maxfs,parker1996iis}:

\begin{theorem}[Forced shifts as a transversal number; standard]\label{thm:transversal}
Call $Y \subseteq H$ \emph{closable} in $P$ if some $\ell \in P$ has
$\shift_\gamma(\ell) = 0$ for all $\gamma \in Y$, and a \emph{core} a set that
is not closable while all its proper subsets are. For any $P$:
(i) $H \setminus X$ is closable if and only if $X$ meets every core; hence
$\tau_P(H)$ is the least size of a set meeting every core;
(ii) for every $\ell \in P$ the turns $\ell$ leaves shifted meet every core;
(iii) every least such set is exactly the set of turns left shifted by some
$\ell \in P$;
(iv) $P \subseteq P'$ implies $\tau_P(H) \ge \tau_{P'}(H)$.
\end{theorem}

\begin{proof}
Closability is hereditary: if $X$ misses a core, $H \setminus X$ contains it,
and a non-closable $H \setminus X$ contains a core. The zero-shift set of
$\ell$ is closed by $\ell$. For a least $X$, a length closing $H \setminus X$
shifts a subset of $X$ meeting every core, hence all of $X$.
\end{proof}

So $\tau_P(H)$ lies between the largest number of disjoint cores and the
least hitting set of the cores found so far, exact once its complement is
closable; the \emph{implicit hitting-set} loop of Parker and
Ryan~\cite{parker1996iis} alternates the two, with
cores from the deletion filter of~\cite{chinneck1991locating}.

\subsection{Whole lengths and size}\label{sec:whole}

\begin{proposition}[Whole lengths and size without a cap]\label{prop:whole-size}
If $P$ is given by the balances and bounds $\ell \ge b \ge 0$, a set closable
in real lengths is closable in whole lengths, and closable in real lengths at
any number of quads greater than $\inf_{P_Y} Q$, and at $\inf_{P_Y} Q$ itself
when the infimum is attained, where $P_Y$ is the set of lengths of $P$ closing
$Y$ and $Q(\ell)$ the number of quads.
\end{proposition}

\begin{proof}
Scale a rational point of $P_Y$ by its denominator (Remark~\ref{rem:exist}).
For a size $q$, take $\ell \in P_Y$ with $0 < Q(\ell) \le q$; $c\ell \in P_Y$
for $c \ge 1$, and $Q(c\ell) = c^2 Q(\ell)$ reaches $q$.
\end{proof}

With a cap on the deviation, $P_Y$ is a polytope and $Q(P_Y)$ an interval: two
exact points on either side of a window of sizes certify that $Y$ closes
within it, while the converse needs the extremes of $Q$ over a polytope. The
cap is a budget shared by all turns: on the body the lengths may move by
$980.7$ edges in all, while closing the $109$ strips of way~2 that the bounds
do not force costs $1{,}680.2$, and the cores become covers of a knapsack
(measured).

On a closed surface where no patch borders itself, an arc lies on a side of
each of the two patches it borders, so its column of $B$ has two entries
$\pm1$: $B$ is the incidence matrix of a \emph{bidirected graph} $G$ whose
nodes are the balances and whose edges are the arcs, an arc on a border giving
a half-edge, the setting of~\cite{heistermann2023bimdf}, with the rank of
Zaslavsky~\cite{zaslavsky1982signed}.
A cycle of $G$ is \emph{balanced} if values $\pm1$ on its edges balance each
of its nodes, and \emph{unbalanced} otherwise.

\begin{lemma}[Handshake]\label{lem:handshake}
Let $x \in \Z^n$ and $C$ be a connected component of $G$. Then $\sum_{r \in C}
(Bx)_r \equiv \sum_{a} x_a \pmod 2$, the second sum running over the border
arcs of $C$. On a closed surface, $\sum_r (Bx)_r$ is even.
\end{lemma}

\begin{proof}
It is $\sum_a x_a \sum_{r \in C} B_{ra}$, and each column sums to $0$ or
$\pm2$, or to $\pm1$ for a border arc.
\end{proof}

So on a closed surface a system $Bx = d$ with $\sum_r d_r$ odd has real
solutions but no whole one; this does not affect the balanced lengths:

\begin{proposition}[Free arcs]\label{prop:free}
Let $G$ be connected and contain an unbalanced cycle. Let $T$ be the edges of a
spanning tree of $G$ and one more edge closing an unbalanced cycle, and $F$ the
other edges. Then $B$ has full row rank, and every $x_F \in \Z^F$ extends to
exactly one $x \in \Z^n$ with $Bx = 0$. The whole balanced lengths thus form a
lattice isomorphic to $\Z^F$, of dimension $n - m$ for $m$ balances, with a
basis of vectors $z_f$, $f \in F$, each supported on $T \cup \{f\}$ with
$z_f(f) = 1$ and entries in $\{0, \pm1, \pm2\}$.
\end{proposition}

\begin{proof}
$T$ exists: the sign of a cycle is a homomorphism on the cycle space over
$\Z/2$, which the fundamental cycles of a spanning tree span, so if they were
all balanced every cycle would be. $T$ has $m$ edges. If $B_T y = 0$, a node
met by a single edge of $T$ forces that edge to zero; pruning such edges
leaves the unbalanced cycle, on which each value fixes the next up to sign,
and going round forces all to zero. So $B_T$ is invertible, $B$ has full row
rank, and $x_F$ has exactly one real extension. For $f \in F$, $T \cup \{f\}$
has a one-dimensional space of balanced values. Pruned, it is either three
paths between two nodes, whose cycles are the unbalanced cycle of $T$, one more
unbalanced and one balanced (the number of unbalanced cycles of such a graph
is even), the balanced one containing $f$, since the cycle of $T$ does not,
and carrying values $\pm1$; or two cycles joined by a path, possibly of length
zero, the cycle through $f$ being balanced and carrying $\pm1$, or unbalanced,
and then $\pm1$ on both cycles and $\pm2$ on the path are balanced. Either way
there is a whole balanced $z_f$ supported on $T \cup \{f\}$ with $z_f(f) =
\pm1$; with signs fixed, $x = \sum_{f \in F} x_f z_f$ is the whole extension.
\end{proof}

With a spanning forest and, in each component with an unbalanced cycle, an
edge closing one or a half-edge, the proposition extends beyond the closed
sphere, where full row rank can fail (on part, $145$ of $148$). On the male
body $G$ is connected with an unbalanced cycle and the lattice of balanced
lengths has dimension $2{,}069 - 1{,}552 = 517$: balances alone never obstruct
whole lengths, and integer elimination ends with a single row of greatest
common divisor $2$. The shifts of the held turns leave lattices of dimension
$363$ to $469$, and in ways 4, 5 and 6 elimination needs steps of Euclid's
algorithm: turn rows can destroy the free-arc structure (measured). The one
integer obstruction found on the body is a rounding gap in deviation, not a
parity: rings that close in real numbers at a mean deviation of $0.399$ need
at least $0.536$ in whole numbers~\cite[\S7]{gardent2026loops}.

\subsection{On the male body}

In Table~\ref{tab:forced}, $H$ is the held turns of a way, or the $111$
strips of the engine's mesh for way~2 leaning less than $0.12$, and $P$ is
$P_{\mathrm{eng}}$, then also the engine's cap $D = 0.474$ on the mean
deviation, then also whole lengths and its window of $10{,}379$ to $10{,}801$
quads. With the bounds alone the implicit hitting-set loop is exact in at most
$66$ seconds per family; with the cap and the window, bounds come from
closable subsets, disjoint cores, solver bounds and whole points checked
exactly.

\begin{table}[t]
\centering
\small
\begin{tabular}{lrrrrr}
\toprule
way & turns & bounds & bounds, cap & cap, size, whole & left shifted by the engine \\
\midrule
0 & 48 & 1 & 1 & 1 & 13 \\
1 & 78 & 5 & 5 & 5--6 & 39 \\
2 & 52 & 1 & 1 & 1 & 16 \\
3 & 86 & 4 & 4 & 5--18 & 44 \\
4 & 84 & 1 & 1 & 1--12 & 49 \\
5 & 160 & 6 & 7--17 & 7--86 & 115 \\
6 & 88 & 1 & 1 & 1--17 & 44 \\
7 & 169 & 9 & 9--26 & 9--90 & 123 \\
\midrule
2, strips of the mesh & 111 & 2 & 3--11 & 3--43 & 81 \\
\bottomrule
\end{tabular}
\caption{Male body: turns of each family that every quantization in $P$ leaves
shifted (Theorems~\ref{thm:forced} and~\ref{thm:transversal}; measured), with
$P = P_{\mathrm{eng}}$, then also the engine's cap (real lengths), then also
whole lengths and the engine's window of sizes. A single number is exact; a
range is certified, below by the solver's proven bound or disjoint cores,
above by a closable subset found within the time limit ($120$ to $1{,}800$
seconds). Last column: turns of the family shifted at the engine's whole
lengths. Turns are counted up to sign.}
\label{tab:forced}
\end{table}

With the bounds alone the forced turns are exactly $1$, $5$, $1$, $4$, $1$,
$6$, $1$ and $9$ (Table~\ref{tab:forced}); in ways 0, 2, 4 and 6 the forced
turn is the wrist; in way~7 nine turns are needed although at most
eight of its $21$ cores are disjoint, so the packing bound is not always
attained. With the cap and the window, at whole lengths, the value is still
exactly $1$ in ways 0 and 2: whole lengths closing every held turn but the
wrist exist at $10{,}439$ and $10{,}436$ quads, at mean deviations $0.442$ and
$0.469$. In way~3 integrality
costs at least one turn ($4$ in real numbers, at least $5$ in whole numbers). Of the
$16$ held turns of way~2 the engine leaves shifted, the bounds force $1$, the
cap, the size and integrality none: the heuristic accounts for $15$; in the
other ways it leaves at least $12$ to $37$ turns more than the best lengths
found. The cap is the engine's only approximately (it refused $0.4737$), and optima
that bring other arcs to zero can break the mesh; keeping open every arc the
engine keeps open, the optimum for way~2 is $4$ (measured).

\section{Choosing the turns}\label{sec:choosing}

Call \RC{} the problem: given the balances $B$, bounds $b$, a set $H$ of
linear forms (the shifts) and a number $k$, are there lengths in $P = \{\ell :
B\ell = 0,\ \ell \ge b\}$ that close at least $k$ forms of $H$? The maximum
feasible subsystem problem is NP-hard and hard to approximate in
general~\cite{amaldi1995complexity}, and minimum IIS covers are strongly
NP-hard already for network flows~\cite{joormann2015flowiis}; the hardness
survives the structure of quantization, a case new as far as we found.

\begin{theorem}[Choosing the turns is NP-hard]\label{thm:nphard}
\RC{} is NP-complete, in real or in whole lengths, already when $B$ is the
incidence matrix of a bidirected graph, the bounds are $0$ or $1$, and the
forms have coefficients $0$ or $1$; it is strongly NP-hard, and stays NP-hard
with a cap on the deviation, which may be loose. Moreover $\tau_P(H)$ can
equal the vertex cover number of any graph, with $|H|$ its number of vertices.
So it cannot be approximated within $\sqrt2 - \varepsilon$ unless
P${}={}$NP~\cite{khot2017independent,dinur2018towards,dinur2018nonoptimally,khot2018grassmann},
nor within $2 - \varepsilon$ if the unique games conjecture
holds~\cite{khot2008vertex}, and the largest closable subset of $H$ cannot be
approximated within $|H|^{1-\varepsilon}$ unless
P${}={}$NP~\cite{hastad1999clique,zuckerman2007linear}.
\end{theorem}

\begin{proof}
Membership: closing a given set of $k$ forms is a linear program, and by
Proposition~\ref{prop:whole-size} real and whole closability agree. Hardness,
from \textup{\textsc{Independent Set}}: for a graph $G$, take for each edge $e
= uv$ four arcs $y_{u,e}$, $y_{v,e}$, $d_e$, $d'_e$ with the balances
\[
  y_{u,e} + y_{v,e} - d_e = 0, \qquad y_{u,e} + y_{v,e} - d'_e = 0, \qquad
  d_e - d'_e = 0 ,
\]
in which every arc appears in exactly two rows with coefficients $\pm1$, the
bounds $d_e, d'_e \ge 1$ and $y \ge 0$, and for each vertex $u$ the form
$\shift_u = \sum_{e \ni u} y_{u,e}$. With $y \ge 0$, $\shift_u = 0$ forces
$y_{u,e} = 0$ at $u$, so closing two adjacent vertices gives $d_e = 0 < 1$.
For an independent set $Y$, every edge has at most one end in $Y$: put $y = 0$
on its side and $1$ on the other, or, with no end in $Y$, $1$ on either side
and $0$ on the other; then $d_e = d'_e = 1$, whole lengths in $P$ closing $Y$.
The closable sets are the independent sets of $G$, and $\tau_P(H) = |V| -
\alpha(G)$ is the vertex cover number. The instance is linear in $G$, with
numbers $0$ and $\pm1$; a large cap constrains nothing.
\end{proof}

The reduction was checked by computer on $40$ random graphs (measured); its
third balance puts every arc in two balances, as in a T-mesh. It is not known
to be a T-mesh: its forms have non-negative coefficients and are not built as
the holonomies of straight walks, so the theorem says that quantization's
algebra allows hard instances, not that a surface produces them.

\begin{conjecture}[Realizable instances]\label{conj:realize}
\RC{} stays NP-hard when $B$, $b$ and $H$ are the balances, the bounds of
separation and the turns of a T-mesh of a closed surface of genus zero.
\end{conjecture}

The evidence is that vertex cover stays NP-hard on planar graphs of small
degree~\cite{garey1976simplified}, and that two turns separated by a
dislocation conflict like the two ends of an edge.

Some parts are easy: without turns and window, whole quantization is a
polynomial matching problem~\cite{edmonds1970matching}; with disjoint known
cores $\tau_P(H)$ is their number; with shift forms of fixed rank the closable
sets are the flats of an arrangement, but on the body the held turns of a way
have full rank.

\begin{remark}[The window of sizes; a sketch]
The window on the number of quads makes even closing nothing hard in the weak
sense. Take a disjoint union of tori, each of two patches whose long and short
sides $a_i$ and $c_i$ have the targets $N v_i$ and $1$, so that $Q = \sum_i 2
a_i c_i$, and a cap that lets the
short sides $c_i$ grow by $T_0$ edges in all. For $N$ large, the window $Q =
2N(\sum_i v_i + T_0)$ is met exactly when $\sum_i v_i(c_i - 1) = T_0$ with
whole $c_i \ge 1$: unbounded subset sum, which is
NP-complete~\cite{lueker1975two}, with numbers exponential in the size of the
instance (the gadgets are of genus one and disconnected). Whether the window
alone is strongly NP-hard, with $Q$ polynomial in the size, is open, as is the
complexity of maximizing $Q$ over a polytope of lengths in reals.
\end{remark}

\section{What is machine-checked}\label{sec:lean}

The folder \texttt{lean} beside this note is a Lean~4~\cite{demoura2021lean4}
project on Mathlib~\cite{mathlib2020} (both \texttt{v4.35.0}). Its $235$
theorem declarations, many of them small lemmas, depend only on Lean's
standard axioms; the count is not evidence, and its \texttt{README.md} maps
each statement to its Lean theorems. Lean checks the algebraic and
combinatorial core of each statement, with the topological facts it uses
(that a loop develops into a given motion, that a curve bounds a disc, that
the loops $\lambda_i$ present the fundamental group, that the quad graph of a
sphere is bipartite) as explicit hypotheses. Each statement is formalized on
its own abstraction (a turn as a list of motions, a region as finite sets with
counts, a system as a pair of matrices): the files share no type of T-mesh
and are not composed into one model of it.

\emph{Checked}, statement by statement. For motions with linear parts in any
ring: the Burgers sum, and that the recursion of the gluing maps of a tree of
cuts, a hypothesis, unfolds into the order of first visit
(Proposition~\ref{prop:burgers}(a)); the cut annulus~(b); exchanges; the
conjugation step of Proposition~\ref{prop:winding}. Over the reals: the four
orders and orientations of a dislocation and straight chains
(Proposition~\ref{prop:dislocation}(i), (ii)); staircases; the positivity of
a non-negative form with a held arc. At whole lengths: the checkerboard rule;
Proposition~\ref{prop:parity} for whole balanced lengths of any sign, given a
2-colouring at positive lengths and the linearity of $\alpha_\gamma$ and
$\shift_\gamma$; the second proof of Remark~\ref{rem:borders}. The
Fox--Burgers formula on the free group with the twisted windings and the
insertion of the relation, the case of a simple turn,
Proposition~\ref{prop:commutator} and Corollary~\ref{cor:checker}.
Lemma~\ref{lem:handshake} on a closed surface; Proposition~\ref{prop:free},
stronger than stated (on the edges $T$ the whole image of $B$ is the lattice
of vectors of even sum); Remarks~\ref{rem:definite} and~\ref{rem:exist}, with
Farkas' lemma over any ordered field and whole multipliers;
Theorem~\ref{thm:forced} and Corollary~\ref{cor:chain}. For fields of any
order: the rotation part of Gauss--Bonnet and its counts on triangle and
hexagonal meshes, the lengths of dislocations and the congruence of
Appendix~\ref{app:orders}.

\emph{Checked in a weaker form}: Proposition~\ref{prop:linear}, only the
additivity of the composed translation, not the integer coefficients nor the
reading as a helix; Lemmas~\ref{lem:gauss} and~\ref{lem:eight}, only the
Euler-characteristic count on the grid refined by two and the totals from the
numbers of convex and concave corners, the regions being discs and the
convexity of their corners hypotheses; the completeness of the walks,
Proposition~2 of~\cite{gardent2026loops}, only in a model: on a closed
surface (opposite sides of a patch of equal length, no two unit edges glued
onto the same one) the station walk of the gluing data is a permutation, the
strips are \emph{defined} as the cycles of the successor it gives on the quads
above the stations, and they are in bijection with the cycles of the real walk
from the half-integers $k + \frac12$, in a one-dimensional model of the
gluing; that patch sides are positive, that crossing a patch keeps the offset
and that the model is the quad mesh are assumptions.
\emph{Not checked}: the topology taken as hypotheses, strips with borders, the
measured statements, and those proved on paper only (see AI assistance). The formalization found no error of sign or
convention; it sharpened the hypotheses of Propositions~\ref{prop:parity},
\ref{prop:free} and~\ref{prop:staircase}, of Remark~\ref{rem:definite} and of
the completeness proposition, stated in their sharpened form.

\section{Open problems and limitations}\label{sec:open}

When is a set of turns closable, in terms of the singularity configuration?
Every certificate is a sum of combs, which says nothing about the
singularities, and not of staircases, caps and winding separatrices alone; a
useful answer lies between, for instance combs on disjoint sets of vertices,
each of the sign of the certificate. Computing $\tau_P(H)$ at a given number of quads needs the extremes of an
indefinite quadratic function over a polytope. Is \RC{} NP-hard for a fixed number of singular vertices? And when
the set of turns depends on the lengths, as the strips of the mesh do, what is
the problem at all? When does a quartic differential with all trajectories finite exist for a
given singularity configuration~\cite{lei2020quartic}? Turns are its discrete
version. Finally, a formal model of a complex
of rectangles with integer sides, its turns and its developing map would turn
the topological hypotheses of the Lean statements into theorems.

\paragraph{Limitations and reproducibility.}
The measured statements come from one T-mesh of one body, exported by an
engine without the constraint rounds of the companion paper, and from a few
bench models with an unrecorded engine version; identities were checked in
floating point unless exact arithmetic is stated, and the homotopies of
Propositions~\ref{prop:swap} and~\ref{prop:winding} were not checked. Forced
counts are relative to the engine's bounds of separation and, approximately,
to its cap, and the intervals depend on solver time limits on a shared
machine. The scripts and logs of the exact experiments are not published: the
measured statements can be checked against the branch \texttt{rd-exact}, an
experimental hook of openclare, but not rerun as they stand. The Lean proofs
replay with \texttt{make proofs NOTE=turns-of-a-t-mesh}. This version of the
note was built from commit \texttt{\enginecommit}.

\section*{AI assistance}
This note, its Lean proofs and the software it describes were made by a human
author working with an AI system, which is not an author. The author chose the
problem and its aims (edge loops that close into rings and, as an artist, few
singular vertices and few spirals rather than perfectly square quads),
directed the work, decided which results to pursue, and judged them through
the measurements and drafts presented to him. Claude Opus 5.5 (Anthropic),
used through Claude Code in October 2026 with several sub-agents, wrote the
changes to the engine and the experiment scripts, ran the exports and exact
computations, searched the literature, stated and proved the propositions,
wrote the Lean formalization and drafted the text. Each bibliography entry says how much of the work was read.
Beyond the cores machine-checked in Lean (Section~\ref{sec:lean}), everything
was written with AI help and checked by AI agents rereading it and testing it
on the data, not by a human expert. In particular, Propositions~\ref{thm:combs},
\ref{prop:pairing}, \ref{prop:collar}, \ref{prop:burgers-any}, \ref{prop:torus}
and~\ref{prop:whole-size}, Lemmas~\ref{lem:comb} and~\ref{lem:gauss-border},
Theorems~\ref{thm:parity-any}, \ref{thm:transversal} and~\ref{thm:nphard}, the
case of Lemma~\ref{lem:handshake} with borders, and the statements on
rotations of order six other than the checked ones have only a paper proof
written with AI help. AI review found errors in Propositions~\ref{thm:combs}
and~\ref{prop:pairing} and Theorem~\ref{thm:nphard}, and in
Proposition~\ref{prop:linear} and the completeness proposition
of~\cite{gardent2026loops}, whose Lean checks are weaker than the statements;
all are corrected here. The note has not been peer reviewed; the author, who
publishes it, is responsible for its content.

\bibliographystyle{plainnat}
\bibliography{refs}

\appendix

\section{Rotations of other orders}\label{app:orders}

On triangle and hexagonal meshes the rotations are powers of $\zeta =
e^{i\pi/3}$, a vertex of valence $k$ has index $6 - k$, and a turn along $u$
comes back translated by $T = \alpha u + \shift\, \zeta u$.
The Burgers sum, the order of first visit and the Fox--Burgers formula do not
use the order of the rotations and hold unchanged, as do
Theorem~\ref{thm:forced}, Corollary~\ref{cor:chain} and
Lemma~\ref{lem:handshake}, which use no geometry. A pair of valences five and seven has $B =
\pm\zeta^{\pm1} d$, and the shifts of two turns differ by $\pm d_1$ or
$\pm(d_1 + d_2)$, $d = d_1 u + d_2 \zeta u$. What changes is parity. When every
rotation is a power $z^{j_i}$ of one rotation $z$, every Fox derivative is
congruent to its exponent sum modulo $1 - z$, and $1 - z^j$ to $j(1 - z)$
modulo $(1 - z)^2$, so in any commutative ring
\[
  \tau(w) \equiv (1 - z) \sum_i \epsilon\Bigl(\frac{\partial w}{\partial
  x_i}\Bigr) j_i\, p_i \pmod{(1 - z)^2}.
\]
For quarter-turns $(1 - z)^2 = -2z$, and this is Corollary~\ref{cor:checker}.
For order six $1 - z$ is a unit: no colouring rule on triangle or hexagonal
meshes. When every singular vertex has even valence the rotations have order
three, $(1 - \omega)^2 = -3\omega$ (compare Lemma~11
of~\cite{izmestiev2013torus}), and on a closed sphere $\alpha \equiv \shift
\pmod 3$: every ring has a length divisible by three, as on subdivided
octahedra, while on subdivided icosahedra rings of $5n$ edges take every
residue (measured); a rule modulo three with borders or handles is a
conjecture. A dislocation of order $n$ has $|B| = 2\sin(\pi/n)\,|d|$.

\end{document}
